\documentclass[11pt]{amsart}

\usepackage[T1]{fontenc}
\usepackage{lmodern}
\usepackage{microtype}
\usepackage{amsmath,amssymb,mathtools}
\usepackage{booktabs}
\usepackage{tabularx}
\usepackage{longtable}
\usepackage{listings}
\usepackage{xcolor}
\usepackage{hyperref}

\lstset{
  basicstyle=\ttfamily\small,
  breaklines=true
}

\hypersetup{
  colorlinks=true,
  linkcolor=blue!55!black,
  citecolor=blue!55!black,
  urlcolor=blue!55!black,
  pdftitle={A conditional trace-73 Cappell--Shaneson skein-lasagna obstruction},
  pdfauthor={Anonymous repository audit}
}

\newtheorem{theorem}{Theorem}[section]
\newtheorem{proposition}[theorem]{Proposition}
\newtheorem{lemma}[theorem]{Lemma}
\newtheorem{corollary}[theorem]{Corollary}
\theoremstyle{definition}
\newtheorem{definition}[theorem]{Definition}
\newtheorem{hypothesis}[theorem]{Hypothesis}
\theoremstyle{remark}
\newtheorem{remark}[theorem]{Remark}
\newtheorem{externalresult}[theorem]{External Result}

\newcommand{\Q}{\mathbb{Q}}
\newcommand{\Z}{\mathbb{Z}}
\newcommand{\Sfour}{S^{4}}
\newcommand{\SN}{\mathcal{S}^{N}_{0}}
\newcommand{\Stwo}{\mathcal{S}^{2}_{0}}
\newcommand{\KhR}{\operatorname{KhR}}
\newcommand{\HH}{\operatorname{HH}_{0}}
\newcommand{\Span}{\operatorname{Span}}
\newcommand{\eps}{\varepsilon}
\newcommand{\Sh}{\operatorname{Sh}}
\newcommand{\im}{\operatorname{im}}
\newcommand{\Bur}{\mathrm{Bur}}
\newcommand{\MWW}{\mathrm{MWW}}
\newcommand{\sha}{\operatorname{SHA256}}
\newcommand{\Open}{\textnormal{\textsc{Open}}}
\newcommand{\Partial}{\textnormal{\textsc{Partial}}}
\newcommand{\Discharged}{\textnormal{\textsc{Discharged}}}
\newcommand{\Unused}{\textnormal{\textsc{Unused}}}
\newcommand{\hash}[2]{\shortstack[l]{\texttt{#1}\\\texttt{#2}}}

\title[A conditional trace-73 obstruction]
{A conditional skein-lasagna obstruction for a trace-73\\
Cappell--Shaneson sphere}

\author{Anonymous repository audit}
\address{Draft prepared from the public trace-73 repository}
\email{Author information to be supplied before submission}
\date{September 3, 2026}

\subjclass[2020]{Primary 57K40, 57R55; Secondary 57K18, 68V20}
\keywords{smooth Poincare conjecture, Cappell--Shaneson sphere,
skein lasagna module, Khovanov homology, formal verification}

\begin{document}

\begin{abstract}
We study the trace-73 Cappell--Shaneson homotopy $4$-sphere associated with
\[
A=\begin{pmatrix}0&269&1240\\0&41&189\\1&0&32\end{pmatrix}.
\]
For an explicit Johnson-generator handle presentation we formulate the
geometric inputs P0, C, S, and P3 needed for a skein-lasagna comparison at
quantum degree $494$, together with the identification $X_J\cong\Sigma_A^0$.
At present P0, C, and S are \Open: the repository certificates check finite
combinatorial models of these inputs, not the actual Cappell--Shaneson
geometry.  An exact finite calculation gives the nonzero integer $2624$ as
the divided cubic of a frozen Burau computation in a fixed collar endpoint
convention; identifying that integer with the MWW divided cubic of an actual
geometric class is part of the open input C.
An Artin--Magnus certificate and the pure-braid Andreadakis theorem
establish a third-order property of the public braid word.
A Lean development formalizes the abstract quotient argument: given interface
data assembling the MWW quotients and four-handle transport, a nonzero
degree-$494$ class would obstruct diffeomorphism with $S^4$.
Those geometric interfaces are not constructed in Lean, so no counterexample
to the smooth four-dimensional Poincar\'e conjecture is claimed.
\end{abstract}

\maketitle
\tableofcontents

\section{Introduction}

The smooth four-dimensional Poincare conjecture asks whether every smooth
closed manifold homotopy equivalent to $S^4$ is diffeomorphic to $S^4$.
Freedman's work proves the topological counterpart: a homotopy 4-sphere is
homeomorphic to $S^4$ \cite{Freedman1982}.  The remaining question is
therefore a smooth classification problem, and any proposed counterexample
must establish two logically separate conclusions:
\begin{enumerate}
\item the candidate is a homotopy 4-sphere;
\item the candidate is not diffeomorphic to the standard smooth sphere.
\end{enumerate}

Cappell and Shaneson constructed a distinguished family of potential
homotopy 4-spheres from torus automorphisms \cite{CappellShaneson1976}.
Their handle structures were studied by Aitchison and Rubinstein
\cite{AitchisonRubinstein1984}.  Many subfamilies were subsequently shown to
be standard, through work including Akbulut \cite{Akbulut2010}, Gompf
\cite{Gompf2010}, Kim--Yamada \cite{KimYamada2017}, and Iwaki
\cite{Iwaki2024}.  This history makes geometric identification especially
important: agreement of a matrix, a relator, or a trace parameter is not a
substitute for identifying a complete framed handle presentation.

Skein lasagna modules extend Khovanov--Rozansky link homology to smooth
four-manifolds \cite{MorrisonWalkerWedrich2019}.  Manolescu and Neithalath
gave a two-handlebody description \cite{ManolescuNeithalath2020}, and
Manolescu, Walker, and Wedrich gave formulas compatible with general handle
decompositions \cite{ManolescuWalkerWedrich2023}.  These results provide an
attractive route to conclusion (ii): a nonzero class in a nonzero bidegree
cannot be carried by a diffeomorphism to the lasagna module of $S^4$, which
is concentrated in bidegree zero.

We investigate the standard-form Cappell--Shaneson parameter
$(41,189,73)$.  A Johnson-generator replacement supplies a finite braid word
and a labelled handle presentation; the paper states the handlebody,
coefficient-comparison, three-handle, four-handle, and
Cappell--Shaneson-identification inputs for that replacement and records,
for each of them, whether it is established or \Open.

\subsection*{Contributions}

The main contributions are as follows.
\begin{enumerate}
\item We record a splitting-preserving Johnson lift and an exact
six-sweep 44-channel word equal to the public 11340-letter word.  This is a
finite P0d fact.  Lemmas P0a--c state the Johnson replacement presentation;
their geometric content (the actual Aitchison--Rubinstein link built from
the actual Johnson homeomorphism, the actual Kirby cancellations, and the
framing identity with the MWW cable) is \Open, and so is the identification
of the closed picture with \(\Sigma_A^0\).
\item The coefficient comparison C and the reversed three-handle picture S
are formulated for the Johnson replacement with their remaining obligations
listed explicitly; both are \Open.
\item A formal Lean core checks the abstract quotient implication; the
status of every geometric instance is summarized in
Section~\ref{sec:status-table} and the replay commands in
Appendix~\ref{sec:reproducibility-extended}.
\end{enumerate}

\section{The trace-73 Cappell--Shaneson candidate}

Let
\begin{equation}\label{eq:A}
A=\begin{pmatrix}
0&269&1240\\
0&41&189\\
1&0&32
\end{pmatrix}.
\end{equation}
This is Iwaki's standard form $X_{41,189,73}$; the triple occurs among the
trace-73 representatives in his table \cite[Table 2]{Iwaki2024}.  Direct
integer arithmetic gives
\begin{equation}\label{eq:determinants}
\det A=1,\qquad \det(A-I)=1.
\end{equation}

For $A\in \operatorname{SL}(3,\Z)$, let $f_A:T^3\to T^3$ be the induced
diffeomorphism, made equal to the identity near a chosen point.  Surgery on
the corresponding circle in the mapping torus gives a
Cappell--Shaneson manifold $\Sigma_A^\varepsilon$, where
$\varepsilon\in\Z/2$ records the framing choice.  The standard criterion says
that $\Sigma_A^\varepsilon$ is a homotopy 4-sphere if and only if
$\det(A-I)=\pm1$; see Iwaki's Proposition 2.1 and the original sources cited
there \cite{Iwaki2024,CappellShaneson1976}.

Consequently, \eqref{eq:determinants} proves that the standard construction
$\Sigma_A^\varepsilon$ is a homotopy 4-sphere.  It does not prove that a
separately generated large Kirby diagram is a diagram of that construction.
That identification is the first and most load-bearing issue below.

Iwaki proves standardness for many infinite families.  The appearance of
$(41,189,73)$ in a representative table is not a standardness or exoticness
theorem for this individual sphere.  Likewise, the trace bound in
Kim--Yamada's work stops below 73 \cite{KimYamada2017}.  These observations
explain the choice of candidate but do not decide its diffeomorphism type.

\section{Precise statements}

We first state the abstract theorem in the form actually proved by the Lean
development.  This prevents geometric intuition from silently strengthening
the result.

Fix a collection of manifold labels, a graded family of rational vector
spaces $G(M,q)$, labels $X$ and $S^4$, and predicates
\(\operatorname{HS}(M)\) and \(\operatorname{Diff}(M,N)\).

\begin{theorem}[Abstract nonstandardness theorem]\label{thm:A}
Suppose there are rational vector spaces $W_0,W_1,W_2,W_3$, an element
$x_0\in W_0$, linear maps
\[
W_0\xrightarrow{q_{01}}W_1\xrightarrow{q_{12}}W_2,
\qquad
\ell_i:W_i\longrightarrow\Q\quad (i=0,1,2),
\]
and linear isomorphisms
\[
T:W_2\xrightarrow{\cong}W_3,
\qquad
F:W_3\xrightarrow{\cong}G(X,494),
\]
such that
\begin{align}
\ell_0(x_0)&=2624,\label{eq:value}\\
\ell_1q_{01}&=\ell_0,\label{eq:q01}\\
\ell_2q_{12}&=\ell_1.\label{eq:q12}
\end{align}
Assume also that every element of $G(S^4,494)$ is zero and that every
diffeomorphism $X\to S^4$ induces a linear isomorphism
\(G(X,494)\cong G(S^4,494)\).  Then $X$ is not diffeomorphic to $S^4$.
\end{theorem}

\begin{proof}
Set $v=F(T(q_{12}(q_{01}(x_0))))$.  If $v=0$, injectivity of $F$ and $T$
implies $q_{12}(q_{01}(x_0))=0$.  Applying $\ell_2$ and using
\eqref{eq:q01}--\eqref{eq:q12} gives $2624=0$, a contradiction.  Thus
$v\ne0$.  If $X$ were diffeomorphic to $S^4$, the induced isomorphism would
send $v$ to an element of $G(S^4,494)$, hence to zero, contradicting
injectivity.
\end{proof}

The Lean theorem corresponding to Theorem~\ref{thm:A} is a conditional
implication from input structures \texttt{ExternalGeometry} and
\texttt{CSExternalGeometry}; no inhabitants are constructed in Lean.
The candidate-specific inputs are stated below and proved in
Sections~\ref{sec:johnson-presentation}--\ref{sec:final-identifications}.

\begin{hypothesis}[Embedded candidate presentation, P0]\label{hyp:P0}
The explicit Johnson-generator replacement is a labelled and framed handle
presentation of $\Sigma_A^0$, with the two framed product cancellations, the
embedded detector ball and the stated cable attachments.
\end{hypothesis}

\begin{remark}[Status of P0]
\Open.  Section~\ref{sec:johnson-presentation} records the Johnson
replacement, and Proposition~\ref{thm:P0discharge} separates what the
finite certificates establish from the obligations that remain.
\end{remark}

\begin{hypothesis}[Coefficient comparison, C]\label{hyp:P1}
The completed MWW coefficient quantum trace admits a grading-preserving,
symmetric-monoidal natural map to the BPW/BHPW weight-one endpoint
representation, natural also for boundary cobordisms supported away from
$B$.  Under this map the point-push action is the public Burau
representation, and the divided cubic row descends through every beta and psi
relation.  Its value on the selected class is the frozen Burau value
$2624$.  The class lies in quantum degree $494$ on the Johnson replacement
collar.
\end{hypothesis}

\begin{remark}[Status of C]
\Open.  See Section~\ref{sec:c-comparison}, in particular
Proposition~\ref{thm:Cdischarge}.
\end{remark}

\begin{hypothesis}[Relative three-handle closure, S]\label{hyp:P2}
The three $3$-handles may be
attached along a complete sphere system disjoint from $B$.  For each sphere
the detector kills the undotted MWW relation and identifies the once-dotted
map with the source term.  Hence the divided detector descends through the
complete three-handle coequalizer.
\end{hypothesis}

\begin{remark}[Status of S]
\Open.  See Section~\ref{sec:s-closure}, in particular
Proposition~\ref{thm:Sdischarge}.
\end{remark}

\begin{hypothesis}[Four-handle closure of the replacement picture, P3]\label{hyp:P3}
After the three $3$-handles of Hypothesis~\ref{hyp:P2} are attached along
the belt spheres of the reversed $1$-handles, the remaining boundary is
$S^3$.  Attaching a $4$-ball along that $S^3$ yields a closed $4$-manifold
$X_J$.  MWW Proposition~3.4 identifies the empty-link lasagna module of this
$W_3$ picture with $\mathcal S^2_0(X_J)$.  MWW Corollary~3.5 gives
$\mathcal S^2_0(S^4)\cong\mathbb Z$ concentrated in bidegree $(0,0)$, so
the quantum-$494$ summand vanishes.  Diffeomorphisms induce
grading-preserving equivalences.  $\det A=\det(A-I)=1$ is the
homotopy-sphere criterion for the standard Cappell--Shaneson construction.
The identification $X_J\cong\Sigma_A^0$ is the constructed Johnson CS
handle picture of Hypothesis~\ref{hyp:P0}, not these determinants.
\end{hypothesis}

\begin{remark}[Status of P3]
The computed part is established: the empty-link Khovanov ranks in quantum
degrees $0$ and $494$, the Euler characteristic of the attaching $4$-ball,
and the grading sum $-44+227+315-4=494$ are computed directly, and the
isomorphisms $\mathcal S^2_0(B^4)\cong\mathcal S^2_0(S^4)$ and
$\mathcal S^2_0(S^4)\cong\mathbb Z$ concentrated in bidegree $(0,0)$ are
\cite[Proposition~3.4 and Corollary~3.5]{ManolescuWalkerWedrich2023}.
The statement that the three $3$-handle attachments along the belt spheres
restore the boundary $S^3$, and the identification $X_J\cong\Sigma_A^0$,
depend on Hypothesis~\ref{hyp:P0} and are therefore \Open.  Independent
replay commands are listed in Appendix~\ref{sec:reproducibility-extended}.
\end{remark}

\begin{theorem}[Conditional trace-73 theorem]\label{thm:joined}
Assume Hypotheses~\ref{hyp:P0}--\ref{hyp:P3} for the Johnson CS handle
picture $X_J$.  Then $X_J$ is identified with $\Sigma_A^0$, Iwaki
Proposition~2.1 makes $\Sigma_A^0$ a homotopy sphere, and if the resulting
geometric data are assembled into the Lean interface
\texttt{ExternalGeometry}, then
\[
\Sigma_A^0\not\cong_{\mathrm{diff}}S^4.
\]
Hypotheses~\ref{hyp:P0}--\ref{hyp:P2} are \Open.
That interface is not constructed in Lean, so no counterexample is claimed.
\end{theorem}

\begin{proof}[Proof of Theorem~\ref{thm:joined}]
Theorem~\ref{thm:A} gives the smooth obstruction once
\texttt{ExternalGeometry} is inhabited.  Hypotheses~\ref{hyp:P0}--\ref{hyp:P2}
supply its geometric fields, and Hypothesis~\ref{hyp:P3} supplies the
four-handle transport and the identification $X_J\cong\Sigma_A^0$.
\end{proof}

\subsection{Status table}\label{sec:status-table}

The table records, for every geometric input, what the replayed
certificates establish and what remains \Open.  It is the claim boundary of
this paper; the claim-boundary checker in the repository reads it.

\begin{center}
\small
\begin{tabularx}{\textwidth}{@{}l l X@{}}
\toprule
Input & Status & Established / remaining \\
\midrule
P0a (handlebody bridge) & \Discharged & Dual-block regular neighbourhoods of the four spines in the barycentric subdivision, explicit elementary-collapse sequences onto the spines, common boundary of genus three, union equal to $T^3$, transport of the complete pair by $S$, and Regina handlebody recognition (Lemma~\ref{lem:P0a}). \\
P0b (two framed cancellations) & \Open & The certificate contains local model cancellation movies in $\mathbb R^3$; the actual Johnson PL homeomorphism $\psi_A$, the actual Aitchison--Rubinstein framed link, and its actual Kirby cancellations with belt spheres, band slides and transported normal fields are not constructed. \\
P0c (MWW cabling framing) & \Open & The $44$ collar strands are a control realization of the public word, not strands cut from the actual cancelled attaching link. \\
P0d (finite word) & \Discharged & The six-sweep Johnson factors reproduce the public $11340$-letter word letter-for-letter, and $\operatorname{lk}(m_2,r_{yz})=0$ in the railroad diagram. \\
C1 (coefficient bimodule) & \Open & The product-ribbon isotopy of the actual cut tangle and the chain-level action squares are not constructed; the certificate checks $44$ model rectangles and $227$ model circles. \\
C2 (statewise cocone) & \Open & The all-cable constant-term statement holds for the endpoint representation; the statewise relations (26)--(27) on every cable state of the actual source, including the local stabilization for $r_{zx}$, are not established. \\
C3 & \Unused & Not used. \\
S (sphere system, hemisphere maps) & \Open & Lemma~\ref{lem:Ssystem} is proved as stated; the actual sphere system in $\partial W_2$, the actual hemisphere movies, and the endpoint factorization for them are not constructed. \\
P2/E7 & \Unused & ``$B$ is fixed'' is not used. \\
P3/E11 & \Open & Depends on Hypothesis~\ref{hyp:P0}. \\
P3/E12 & \Discharged & Computed vanishing of the empty-link degree-$494$ summand on the standard $S^4$; MWW Proposition~3.4 and Corollary~3.5 are cited. \\
P3/E13 & \Open & The identification $X_J\cong\Sigma_A^0$ depends on Hypothesis~\ref{hyp:P0}. \\
Finite detector & \Discharged & $2624$ is the exact value of the frozen Burau computation in the single collar endpoint convention. \\
Lean \texttt{ExternalGeometry} & \Open & No inhabitant is constructed. \\
\bottomrule
\end{tabularx}
\end{center}

\section{Published results used, with their complete scope}\label{sec:published-results}

This section records the mathematical content imported from the literature.
The wording is a faithful restatement rather than a quotation.  In each case
the hypotheses are at least those in the cited source, and the conclusion is
not strengthened.  Candidate-specific identifications are supplied by
Hypotheses~\ref{hyp:P0}--\ref{hyp:P3} or remain as uninhabited Lean fields;
they are not added to the external statements themselves.

\subsection{One- and two-handle formulas}

\begin{externalresult}[MWW Theorem~4.7]\label{prop:mww-4.7}
Let $k$ be a field, let $W_1=\natural^m(S^1\times B^3)$, and let
$L\subset\partial W_1$ be a null-homologous framed link.  Suppose that $L$
meets the belt sphere of the $i$th one-handle transversely in $2p_i$ points.
Cutting along these belt spheres gives a tangle $R$ in the complement of
paired balls in $S^3$.  Then the gluing map induces an isomorphism from the
direct sum, over tangles $T_i$ with the prescribed boundary, of
\[
 \KhR_N\!\left(R\cup\bigsqcup_i(T_i\sqcup\overline{T_i});k\right)
 \left\{\Bigl(\sum_i p_i\Bigr)(N-1)\right\}
\]
modulo the two gluing actions of the $3$-ball tangle categories, to
$\SN(W_1;L;k)$.  The source states that a global grading shift is omitted
from this displayed presentation.
\end{externalresult}

\paragraph{Source and scope.}
This is External Result~E1, reproduced with the field, null-homology,
transversality, boundary, and grading-shift hypotheses of
Manolescu--Walker--Wedrich \cite[Theorem~4.7]{ManolescuWalkerWedrich2023}.
Its candidate-specific application is the Johnson replacement collar of
Hypothesis~\ref{hyp:P1}.

\begin{externalresult}[MWW Definition~3.1 and Theorem~3.2]\label{prop:mww-3.2}
Let $W$ be a smooth oriented four-manifold, let $L\subset\partial W$ be a
framed link, and let $W'$ be obtained by attaching two-handles along a framed
link $K\subset\partial W$ disjoint from $L$.  For every relative class
$(\alpha,\eta)\in H_2^L(W';\Z)$, form the direct sum over $r\in\mathbb N^n$ of the
cabled summands
\[
 \SN\!\left(W;K(r-\alpha^-,r+\alpha^+)\cup L,\eta^r\right)
 \{(1-N)(2|r|+|\alpha|)\}.
\]
Quotient by all signed-colour-preserving braid relations
$\beta_i(b)v\sim v$ and all stabilization relations
$\psi_i^{[d]}(v)\sim0$ for $d<N-1$ and
$\psi_i^{[N-1]}(v)\sim v$.  Attaching the corresponding positive and
negative parallel core disks defines an isomorphism from this cabled quotient
to $\SN(W';L;(\alpha,\eta))$.
\end{externalresult}

\paragraph{Source and scope.}
This is exactly the scope of MWW Definition~3.1 and Theorem~3.2
\cite[Definition~3.1 and Theorem~3.2]{ManolescuWalkerWedrich2023}.
The anti-parallel braid relations are retained; we do not replace their
braid groups by symmetric groups.

\subsection{Three- and four-handle formulas}

\begin{externalresult}[MWW Theorem~3.7]\label{prop:mww-3.7}
Let $W$ be a smooth oriented four-manifold with boundary $Y$, let
$S\subset Y$ be an embedded $2$-sphere disjoint from a framed link
$L\subset Y$, and let $W'$ be obtained by attaching a three-handle along
$S$.  Let $J$ be an oriented equator of $S$, and let $\Delta_+$ and
$\Delta_-$ be the two hemispheres, pushed into $I\times Y$ and adjoined to
$I\times L$.  For a fixed outgoing relative class $\alpha'$, take the direct
sum over all incoming relative classes that reduce to $\alpha'$ modulo
$[S]$.  Then the three-handle map to
$\SN(W';L';\alpha')$ is surjective, its kernel is the image of the
difference of the two hemisphere maps, and the target is their coequalizer.
\end{externalresult}

\paragraph{Source and scope.}
This is MWW Theorem~3.7 \cite[Theorem~3.7]{ManolescuWalkerWedrich2023}.
The direct sum over relative classes is part of the theorem.

\begin{externalresult}[MWW Theorem~3.10]\label{prop:mww-3.10}
Suppose
$W_1\subset W_2\subset W_3\subset W_4$ is a fixed handle decomposition in
which $W_1=\natural^m(S^1\times B^3)$, $W_2$ is obtained by attaching
two-handles along $K$, $W_3$ by attaching three-handles along
$S_1,\ldots,S_p$, and $W_4$ by attaching four-handles.  Let
$L\subset\partial W_4$ be a framed link, represent the spheres by genus-zero
surfaces in $\partial W_1$ completed by two-handle cores, and fix a relative
class $\alpha'$.  Then $\SN(W_4;L;\alpha')$ is the quotient of the
direct sum of the cabled modules over all lifts of $\alpha'$ by, for each
three-handle sphere, the undotted through $(N-2)$-dotted relations equal to
zero and the $(N-1)$-dotted relation equal to the identity.
\end{externalresult}

\paragraph{Source and scope.}
This is MWW Theorem~3.10 with its handle-decomposition, surface-realization,
link, and relative-class hypotheses
\cite[Theorem~3.10]{ManolescuWalkerWedrich2023}.
For $N=2$ it yields one undotted relation and one once-dotted-minus-identity
relation for each sphere.

\begin{externalresult}[MWW Proposition~3.4 and Corollary~3.5]\label{prop:mww-3.4}
For the empty boundary link, inclusion of a four-manifold into the result of a
three-handle attachment induces a surjection on $\SN$, while inclusion
into the result of a four-handle attachment induces an isomorphism.  Over the
integral coefficient convention of the source,
$\SN(S^4)\cong\Z$ and is concentrated in bidegree $(0,0)$.
\end{externalresult}

\paragraph{Source and scope.}
This is MWW Proposition~3.4 and Corollary~3.5
\cite[Proposition~3.4 and Corollary~3.5]{ManolescuWalkerWedrich2023}.
The source uses the two-handlebody identification of
Manolescu--Neithalath \cite{ManolescuNeithalath2020} in the handle formula.
The source corollary is integral.  The repository does not prove
Proposition~3.4 or Corollary~3.5; Hypothesis~\ref{hyp:P3} cites them and
separates the computed empty-link ranks from those citations.

\begin{externalresult}[Morrison--Walker--Wedrich Definition~5.4 and Example~5.6]\label{prop:mww-B4}
For a smooth oriented four-manifold $W$ and a framed oriented link
$L\subset\partial W$, the group $\SN(W;L)$ is the bigraded abelian group
generated by labelled lasagna fillings modulo multilinearity, replacement of
an input ball by a filling with the same $\KhR_N$ evaluation, and isotopy
relative to the boundary.  If $W=B^4$ and $L\subset S^3=\partial B^4$, the
evaluation of fillings induces an isomorphism
\begin{equation}\label{eq:B4-evaluation}
\SN(B^4;L)\xrightarrow{\cong}\KhR_N(S^3,L).
\end{equation}
\end{externalresult}

\paragraph{Source and scope.}
This is the invariant definition and $4$-ball evaluation in
Morrison--Walker--Wedrich
\cite[Definition~5.4 and Example~5.6]{MorrisonWalkerWedrich2019}.
The source states the integral bigraded group.  Lean
\texttt{S4ReductionData} packages only the empty-link control on
\texttt{EmptyKhQ}; it does not inhabit candidate \texttt{ExternalGeometry}.

\subsection{Replacement of the three attaching spheres}

\begin{externalresult}[Horvat--Jab\l onowski Theorem~5.3]\label{prop:hj-theorem53}
Let $W$ be a smooth compact connected four-manifold with nonempty connected
boundary, equipped with a fixed handle decomposition without four-handles.
Suppose the decomposition has $k$ three-handles and
\[
 \partial W_2\cong M'\#\bigl(\#_k(S^1\times S^2)\bigr),
\]
where $M'$ is irreducible.  For pairwise-disjoint smoothly embedded
$2$-spheres $\Sigma_1,\ldots,\Sigma_k\subset\partial W_2$, the following are
equivalent: they may be isotoped, up to permutation and three--three handle
slides, to the actual attaching spheres; their exterior is connected; and
their homology classes form a $\Z$-basis of
$H_2^{\mathrm{sph}}(\partial W_2)$.  When these conditions hold, simultaneous
surgery on the spheres produces a closed $3$-manifold diffeomorphic to $M'$.
\end{externalresult}

\paragraph{Source and scope.}
This is the full statement of Theorem~5.3 in
Horvat--Jab\l onowski \cite[Theorem~5.3]{HorvatJablonowski2025}.
The boundary decomposition, irreducibility, number of spheres, pairwise
disjointness, and fixed-handle-decomposition hypotheses are all retained.
Closed-manifold Theorem~5.3 is not used to conclude that the detector ball
$B$ is fixed.  Hypothesis~\ref{hyp:P2} uses it only with
Lemma~\ref{lem:Ssystem} for kernel invariance.  The source (arXiv v3) also
contains Lemma~5.5 (spotted-ball slide lemma) and Lemma~5.7 (relative
uniqueness of complete sphere systems), through which Theorem~5.3 is
proved; neither lemma is invoked directly in the Johnson replacement
argument, and the relative form of Lemma~5.7 is not used to fix $B$.

\subsection{Trace, functoriality, and the Cappell--Shaneson criterion}

\begin{externalresult}[BPW Proposition~3.20 and Theorem~3.21]\label{prop:bpw-trace}
For a locally pregraded endobicategory $(\mathcal C,\Sigma)$ with left duals,
the quantum horizontal trace is universal among $\Sigma$-twisted quantum
preshadows; if right duals also exist, it is a quantum shadow.  On the
$2$-category of locally pregraded endobicategories with duals and
degree-preserving structure maps, the quantum horizontal trace is a strict
$2$-functor.
\end{externalresult}

\paragraph{Source and scope.}
This combines, without enlarging either conclusion, Beliakova--Putyra--Wehrli
Proposition~3.20 and Theorem~3.21
\cite[Proposition~3.20 and Theorem~3.21]{BeliakovaPutyraWehrli2019}.

\begin{externalresult}[BHPW Theorem~4.6]\label{prop:bhpw}
For an oriented tangle $T$, the homotopy type of the Blanchet--Khovanov
bracket is an invariant of $T$ and is strictly functorial with respect to
ordinary tangle cobordisms.  Under the stated BHPW equivalence from the foam
category to the oriented Bar--Natan category, its image is isomorphic to the
ordinary Khovanov bracket.
\end{externalresult}

\paragraph{Source and scope.}
This is the strictification result in Beliakova--Hogancamp--Putyra--Wehrli
\cite[Theorem~4.6]{BeliakovaHogancampPutyraWehrli2023}.  The immediately
following remark in the source labels the extension to knotted webs and
four-dimensional singular foams as conjectural.  Those objects are excluded
from Hypothesis~\ref{hyp:P1}.

\begin{externalresult}[Iwaki Proposition~2.1]\label{prop:cs-criterion}
Let $A\in\mathrm{SL}(3,\Z)$ and let $\Sigma_A^\varepsilon$ be obtained from the
mapping torus of the induced automorphism of $T^3$ by surgery on the canonical
section circle with either of the two framing parities.  Then
$\Sigma_A^\varepsilon$ is a homotopy $4$-sphere if and only if
$\det(A-I)=\pm1$.
\end{externalresult}

\paragraph{Source and scope.}
This is Iwaki Proposition~2.1, citing the original Cappell--Shaneson
construction \cite[Proposition~2.1]{Iwaki2024,CappellShaneson1976}.
The identification of the Johnson picture $X_J$ with $\Sigma_A^0$ is
Hypothesis~\ref{hyp:P0} together with P3/E13, not part of Iwaki's theorem.

\subsection{Standard topology results represented in Lean}

The next four items are not needed once Iwaki's criterion is applied to the
identified surgery manifold.  They state the longer topology route represented
in \path{Smooth4PC/T73CSTopology.lean}.  Lean stores their candidate-specific
application as fields of \texttt{CSTopologyData}; no inhabitant of that
structure is constructed.

\begin{externalresult}[Seifert--van Kampen theorem]\label{prop:van-kampen}
Let $X=U\cup V$, where $U$, $V$, and $U\cap V$ are open and path connected,
and choose a basepoint in $U\cap V$.  Then the diagram induced by inclusion
\[
\pi_1(U\cap V)\longrightarrow\pi_1(U),
\qquad
\pi_1(U\cap V)\longrightarrow\pi_1(V)
\]
is a pushout diagram of groups: $\pi_1(X)$ is the free product of
$\pi_1(U)$ and $\pi_1(V)$ modulo the normal closure of the two images of each
element of $\pi_1(U\cap V)$.
\end{externalresult}

\paragraph{Source and scope.}
This is the standard based, path-connected form of van Kampen
\cite[Section~1.2]{Hatcher2002}.  Applying it to Cappell--Shaneson surgery
still requires the geometric identification of the cover and of the map
$A-I$.

\begin{externalresult}[Wang and Mayer--Vietoris exact sequences]\label{prop:wang-mv}
For a fibration $F\to E\to S^1$ with monodromy $f:F\to F$, there is a long
exact Wang sequence whose relevant portion is
\[
\cdots\to H_k(F)\xrightarrow{\,I-f_*\,}H_k(F)
\longrightarrow H_k(E)\longrightarrow
H_{k-1}(F)\xrightarrow{\,I-f_*\,}H_{k-1}(F)\to\cdots.
\]
For an excisive decomposition $X=A\cup B$, singular homology admits the long
exact Mayer--Vietoris sequence
\[
\cdots\to H_k(A\cap B)\to H_k(A)\oplus H_k(B)
\to H_k(X)\to H_{k-1}(A\cap B)\to\cdots,
\]
with the standard inclusion-induced maps and connecting homomorphism.
\end{externalresult}

\paragraph{Source and scope.}
These are the ordinary singular-homology exact sequences
\cite[Sections~2.2 and 4.2]{Hatcher2002}.  The candidate-specific claim that
their maps are the matrices in \eqref{eq:A-minus-I} and \eqref{eq:H2-map}
is not part of the general theorem.

\begin{externalresult}[Integral Poincar\'e duality]\label{prop:PD}
If $M$ is a closed connected oriented topological $n$-manifold with
fundamental class $[M]\in H_n(M;\Z)$, cap product with $[M]$ gives an
isomorphism
\[
H^k(M;\Z)\xrightarrow{\cong}H_{n-k}(M;\Z)
\]
for every $k$.
\end{externalresult}

\paragraph{Source and scope.}
This is the integral oriented form of Poincar\'e duality
\cite[Section~3.3]{Hatcher2002}.  Closedness, connectedness, orientation, and
integral coefficients are retained.

\begin{externalresult}[Hurewicz and Whitehead promotion]\label{prop:hurewicz-whitehead}
If a path-connected space $X$ is $(n-1)$-connected for $n\ge2$, then
$H_i(X;\Z)=0$ for $i<n$ and the Hurewicz map
$\pi_n(X)\to H_n(X;\Z)$ is an isomorphism.  If a map between simply connected
CW complexes induces an isomorphism on all integral homology groups, then it
is a homotopy equivalence.
\end{externalresult}

\paragraph{Source and scope.}
The first sentence is the Hurewicz theorem and the second is the homological
form of Whitehead's theorem for simply connected CW complexes
\cite[Sections~4.1--4.2]{Hatcher2002}.  Applied successively, they show that a
simply connected CW complex with the integral homology of $S^4$ is homotopy
equivalent to $S^4$.  The premise that the candidate has that topology is
external to the abstract theorems.

\begin{remark}[Why the short route is used]\label{rem:topology-tools}
The homotopy-sphere conclusion for $\Sigma_A^0$ uses Iwaki
Proposition~2.1 rather than silently rebuilding the Cappell--Shaneson
calculation from these four results.  The longer route remains as Lean
fields of \texttt{CSTopologyData}; those fields are uninhabited.
\end{remark}


\section{The finite computational layer}

The public finite layer has three independent parts: lattice arithmetic, a
chosen-sphere coordinate determinant, and a truncated Burau calculation.

\subsection{Integral lattice calculations}

Besides \eqref{eq:determinants}, the repository records the proposed sphere
coordinate columns
\begin{equation}\label{eq:spherecols}
C=\begin{pmatrix}
-1311&-189&41\\
8608&1241&-269\\
-1&0&1
\end{pmatrix},
\qquad \det C=1.
\end{equation}
Thus the three recorded coordinate vectors form a basis of $\Z^3$.  This is
only the last arithmetic step in a geometric-basis argument.  It does not
construct embedded spheres, prove disjointness, or identify $\Z^3$ with
$H_2^{\mathrm{sph}}(\partial W_2)$.

The Lean files also exhibit integral inverses for $A-I$ and for the induced
map on $H_2(T^3;\Z)$.  These computations eliminate the candidate-specific
lattice algebra from the homotopy-sphere branch.  Van Kampen,
Wang--Mayer--Vietoris, Poincare duality, and Hurewicz--Whitehead arguments
remain supplied through a topology structure rather than formalized as
four-manifold topology.

\subsection{The public Burau computation}

Let $\varepsilon=t-1$.  From $252$ primitive crossing rows one reconstructs
an $11{,}340$-letter Artin word on $44$ strands, cables it to $88$ strands,
and applies the unreduced Burau representation over
$\Z[\varepsilon]/(\varepsilon^7)$ with $t=q^{-2}$, $q=1+h$, and
$\varepsilon=(1+h)^{-2}-1$.

Denote the resulting scalar by $\delta_3^{\Bur}$.  The independently
reproduced public result is
\begin{equation}\label{eq:burau}
\delta_3=2624=(-2)^3\cdot(-328).
\end{equation}
An earlier draft mixed two endpoint coordinate tables and reported
$-59072$; with the collar endpoints of Section~\ref{sec:exact-algebra} the
value is $2624$, which is still nonzero.\footnote{Independent replay:
\texttt{scripts/recompute\_t73\_delta3.py --check}.}

\begin{proposition}[Exact scope of the computation]\label{prop:scope}
Equation~\eqref{eq:burau} concerns the truncated Burau model with the collar
endpoints of \eqref{eq:u} and \eqref{eq:ell}.  Identifying the result with
the MWW divided cubic is the content of Hypothesis~\ref{hyp:P1}, which is
\Open.
\end{proposition}

\begin{proof}
The Burau calculation supplies an integer; Hypothesis~\ref{hyp:P1} and
Section~\ref{sec:c-comparison} would identify it with the MWW detector on
the Johnson replacement collar.  Until that hypothesis is established,
$2624$ is only the value of the frozen Burau computation.
\end{proof}

Thus the finite recomputation and the categorical comparison have distinct
roles: the former supplies the integer, while the latter makes it a
four-manifold invariant.

\section{Exact finite algebra and the public detector}\label{sec:exact-algebra}

This section makes the finite layer of Section~4 self-contained.  The
geometric binding of the same data to the Johnson replacement collar is
Hypothesis~\ref{hyp:P1}, which is \Open; the integer $2624$ itself is
independent of that binding and is the value of the frozen Burau
computation.

\subsection{Erratum to the 2 September public draft}

The 2 September public draft reported $D_3(v_T)=-59072$.  That value mixed
two endpoint coordinate systems: the vector $u$ used the THXY endpoint
indices, whereas the covector $\ell$ used the collar indices that also label
the cabled Artin word.  The two tables order the $m_2$ wickets in opposite
directions.  In the collar table, the same physical cup has endpoints $2$
and $87$, so
\[
  u=e_2-e_{87},\qquad \ell=e_{87}^*-e_2^*.
\]
With both objects in that one table, the exact calculation gives
$[\varepsilon^3]\ell(\rho(W)-I)u=-328$ and hence the cubic coefficient is
$(-2)^3(-328)=2624$.  The nonvanishing conclusion is unchanged.
Placing both $u$ and $\ell$ in the THXY coordinates while leaving the word
in collar coordinates gives $-2496$; this is not a coordinate change of the
whole pairing, because the operator is not transported, and it is not a
second geometric class (Appendix~\ref{app:legality}).  The endpoint
transport program reproduces the withdrawn numbers as illegal partial
transports: THXY cup with collar cap, $-59072$; THXY cup and cap with the
untransported word, $-2496$; frozen cut-slot cup with collar cap,
$-55104$; cut-slot cup and cap with the untransported word, $-115456$.
Simultaneous transport of the word, the cup and the cap gives $2624$ in
every coordinate system, and the derived constant terms are exactly
\eqref{eq:u} and \eqref{eq:ell}.  The value $2624$ belongs to the
currently frozen collar data; if the Johnson collar word changes, the
numerical value must be recomputed.

\subsection{Integral inverses and the long topology route}

Besides \eqref{eq:determinants}, define
\begin{equation}\label{eq:A-minus-I}
A-I=
\begin{pmatrix}
-1&269&1240\\
0&40&189\\
1&0&31
\end{pmatrix}.
\end{equation}
Expansion along the first row gives $\det(A-I)=1$.  An integral inverse is
\begin{equation}\label{eq:B-inverse}
(A-I)^{-1}_{\Z}=
\begin{pmatrix}
1240&-8339&1241\\
189&-1271&189\\
-40&269&-40
\end{pmatrix}.
\end{equation}
Since $\det A=1$, direct inversion gives
\begin{equation}\label{eq:H2-map}
C=(A^{-1})^T-I=
\begin{pmatrix}
1311&189&-41\\
-8608&-1241&269\\
1&0&-1
\end{pmatrix},
\end{equation}
with integral inverse
\begin{equation}\label{eq:H2-inverse}
C^{-1}_{\Z}=
\begin{pmatrix}
-1241&-189&40\\
8339&1270&-269\\
-1241&-189&39
\end{pmatrix}.
\end{equation}
These matrices satisfy $\det A=\det(A-I)=1$.  The sphere coordinate matrix in
\eqref{eq:spherecols} is obtained by negating the columns of
\eqref{eq:H2-map} and also has determinant $1$.

\begin{remark}[Long topology route]\label{rem:long-topology-route}
Surjectivity of $A-I$ kills the natural coinvariant presentation of the
fundamental group after surgery; bijectivity of $A-I$ and $(A^{-1})^T-I$
kills the intermediate homology in the Wang and Mayer--Vietoris sequences.
Poincar\'e duality then gives the homology profile of $S^4$.  Hurewicz and
Whitehead promote simple connectivity plus that homology profile to a
homotopy equivalence with $S^4$.  This paper uses Iwaki
Proposition~2.1 on the identified surgery manifold instead.
\end{remark}

\begin{corollary}[Homotopy-sphere conclusion]
\label{cor:homotopy-sphere}
The standard Cappell--Shaneson manifold $\Sigma_A^0$ is a homotopy
$4$-sphere.  If the Johnson CS handle picture $X_J$ is identified with
$\Sigma_A^0$ (P3/E13, \Open), the same holds for $X_J$.
\end{corollary}

\begin{proof}
By \eqref{eq:determinants}, $A\in\mathrm{SL}(3,\Z)$ and $\det(A-I)=1$, so
Iwaki Proposition~2.1 \cite{Iwaki2024} applies to $\Sigma_A^0$.  The
identification $X_J\cong\Sigma_A^0$ depends on Hypothesis~\ref{hyp:P0}.
\end{proof}

\subsection{Coefficient traces and the selected Hattori class}

\begin{definition}[Coefficient trace]\label{def:coefficient-HH0}
For a small $\Q$-linear category $\mathcal C$ and coefficient bimodule $M$,
the coefficient zeroth Hochschild homology is
\[
\HH(\mathcal C;M)=
\left(\bigoplus_{T}M(T,T)\right)
\Big/
\Span_{\Q}\{f\cdot m-m\cdot f\},
\]
where $f:S\to T$ and $m\in M(T,S)$ range over every cyclically composable
pair.
\end{definition}

\begin{lemma}[Descent through the coefficient trace]\label{lem:HH0-descent}
Let $r_T:M(T,T)\to V$ satisfy $r_S(f\cdot m)=r_T(m\cdot f)$ for all
cyclically composable $f$ and $m$.  Then there is a unique linear map
$\bar r:\HH(\mathcal C;M)\to V$ whose restriction to $M(T,T)$ is $r_T$.
\end{lemma}

\begin{proof}
The sum of the $r_T$ kills every relation generator $f\cdot m-m\cdot f$ and
therefore factors uniquely through the quotient.
\end{proof}

For the Johnson replacement collar, Lemma~\ref{lem:C1}, once established,
supplies the actual coefficient equivalence \eqref{eq:HattoriActual}; at
present it is \Open.  The selected class and
divided detector of Section~7 are the geometric instances of the following
pattern.

\subsection{The endpoint module and the Burau action}

\begin{definition}[Endpoint module]\label{def:endpoint-module}
Let $R_\eps=\Z[\eps]/(\eps^7)$, $t=1+\eps$, and
$E_\eps=R_\eps^{88}=\bigoplus_{i=0}^{87}R_\eps e_i$.  One may identify this
free module with the weight-$86$ subspace of the $88$-fold tensor power of
the two-dimensional fundamental module.
\end{definition}

\begin{definition}[Unreduced Burau representation]\label{def:burau}
For the Artin generator $\sigma_i$ of $B_{88}$, with $1\le i\le87$, define
$\rho_t(\sigma_i)$ to be the identity off the span of
$e_{i-1},e_i$ and to have the block
\[
\rho_t(\sigma_i)\big|_{\langle e_{i-1},e_i\rangle}
=
\begin{pmatrix}1-t&t\\1&0\end{pmatrix}
=
\begin{pmatrix}-\eps&1+\eps\\1&0\end{pmatrix}.
\]
Its inverse has block
\[
\rho_t(\sigma_i^{-1})\big|_{\langle e_{i-1},e_i\rangle}
=
\begin{pmatrix}0&1\\t^{-1}&1-t^{-1}\end{pmatrix},
\]
where $t^{-1}=1-\eps+\eps^2-\eps^3+\eps^4-\eps^5+\eps^6$ in $R_\eps$.
For a word $w=\sigma_{i_1}^{s_1}\cdots\sigma_{i_m}^{s_m}$ we use the left
action $\rho_t(w)v=\rho_t(\sigma_{i_1}^{s_1})\cdots
\rho_t(\sigma_{i_m}^{s_m})v$.
\end{definition}

\subsection{From primitive crossings to the cabled word}

The public input stores $252$ primitive crossing rows.  For $i<j$ define the
signed pure-row words
\begin{align}
L_{ij}^{s}
 &=\sigma_{j-1}\cdots\sigma_{i+1}
   \sigma_i^{s}\sigma_i^{s}
   \sigma_{i+1}^{-1}\cdots\sigma_{j-1}^{-1},
\label{eq:L-row}\\
R_{ij}^{s}
 &=\sigma_i\cdots\sigma_{j-2}
   \sigma_{j-1}^{s}\sigma_{j-1}^{s}
   \sigma_{j-2}^{-1}\cdots\sigma_i^{-1}.
\label{eq:R-row}
\end{align}
The $44$-strand word $B_{44}$ is obtained by taking the $R$ rows on segment
$6$ in increasing horizontal coordinate, followed by the $L$ rows on segment
$2$ in decreasing horizontal coordinate.  The cabling homomorphism is
\begin{equation}\label{eq:cabling}
c(\sigma_i)=\sigma_{2i}\sigma_{2i+1}\sigma_{2i-1}\sigma_{2i},
\qquad
c(\sigma_i^{-1})=c(\sigma_i)^{-1},
\end{equation}
and the $88$-strand word is $B_{88}=c(B_{44})$.

\begin{proposition}[Word identities]\label{prop:word-identities}
The exact reconstruction gives
$|B_{44}|=11340$ and $|B_{88}|=45360$.  The Johnson reconstruction verifier
recovers the same $11340$-letter word letter-for-letter.
\end{proposition}

\subsection{The detector and the exact cubic}

Let $R_h=\Q[h]/(h^7)$ and substitute $q=1+h$, $t=q^{-2}$, $\eps=t-1$.
Under the endpoint normalization of the public input, the composite row is
\begin{equation}\label{eq:ell}
\ell=e_{87}^{*}-e_2^{*},
\end{equation}
and the selected shadow vector is
\begin{equation}\label{eq:u}
u=e_2-e_{87}.
\end{equation}
Both formulas use the collar-coordinate endpoint table fixed in the public
input.

\begin{definition}[The divided cubic detector]\label{def:detector}
For a quantum trace class $x$, define
\[
\mathcal D_h(x)=
\widehat C(h)\bigl(\rho_h(W)-I\bigr)P_{86}\Sh_h(x),
\]
where $\Sh_h$ is the quantum trace shadow, $P_{86}$ projects onto the
through-$86$ top cell, and $\widehat C(h)$ is the normalized endpoint cap
row.  When the image lies in $h^3R_h$, the divided cubic detector is
\begin{equation}\label{eq:D3}
D_3(\bar x)=\left.\frac{\mathcal D_h(x)}{h^3}\right|_{h=0}.
\end{equation}
\end{definition}

\begin{proposition}[Selected value]\label{prop:D3-vT}
In $E_\eps$, the first nonzero coefficient of $(\rho_t(B_{88})-I)u$ occurs in
degree $3$.  Contracting with \eqref{eq:ell} gives
\begin{equation}\label{eq:ell-eps-series}
\ell(\rho_t(B_{88})-I)u
=-328\eps^3+14596\eps^4-410246\eps^5+9595271\eps^6.
\end{equation}
Hence, with $[h]\eps=-2$,
\begin{equation}\label{eq:D3-vT-value}
[h^3]\,\ell(\rho_h(B_{88})-I)u=2624\ne0,
\end{equation}
which is $D_3([v_T])$ once Hypothesis~\ref{hyp:P1} identifies $[v_T]$ with
an actual coefficient-trace class.
The complete degree-three vector is printed in Appendix~\ref{app:vector}.
\end{proposition}

\begin{proof}
Traverse $B_{88}$ from right to left, reduce modulo $\eps^7$, and contract
with \eqref{eq:ell}.  The resulting series is
\eqref{eq:ell-eps-series}, so the cubic coefficient is $(-2)^3(-328)=2624$.
\end{proof}

\begin{lemma}[Parameter expansion]\label{lem:eps-of-h}
In $\Z[h]/(h^7)$,
$\eps=(1+h)^{-2}-1=-2h+3h^2-4h^3+5h^4-6h^5+7h^6$.
\end{lemma}

\begin{lemma}[Commutator order]\label{lem:commutator-order}
If $P=I+O(h^a)$ and $Q=I+O(h^b)$ are invertible matrices over an
$h$-adically filtered ring, then $PQP^{-1}Q^{-1}=I+O(h^{a+b})$.  If every
generator of a pure-braid subgroup acts as $I+O(h)$, then its $r$th
lower-central subgroup acts as $I+O(h^r)$.
\end{lemma}

\begin{remark}[Relative class]\label{rem:vT-xi}
The relative class $\xi=\eta_R[T_0]-s_{\mathrm{inv}}\eta_R[T_1]$ satisfies
$D_3(\xi)=0$ because the detector already contains one factor
$\rho_h(W)-I$.  The plain shadow coefficient $[h^3]\ell\Sh_h(\xi)$ equals
$2624$, but that plain shadow is not $D_3(\xi)$.
\end{remark}

\subsection{The grading calculation}

\begin{proposition}[Quantum degree]\label{prop:degree-494}
For the Johnson replacement collar, the selected raw class has quantum
degree $494$.
\end{proposition}

\begin{proof}
There are four contributions:
\begin{align}
d_1&=-44 &&\text{(remove the $44$ Hom normalizations for $N=2$)},\\
d_2&=+227 &&\text{($227$ distinct Frobenius labels)},\\
d_3&=+315 &&\text{($44$ $y$-pairs and $271=44+227$ $z$-pairs)},\\
d_4&=-4 &&\text{(cabled shift for $N=2$, $\alpha=0$, $|r|=2$)}.
\end{align}
Adding gives $-44+227+315-4=494$.  Appendix~\ref{app:grading} collects the
same table without abbreviation.
\end{proof}

\subsection{Quotient algebra and the abstract descent chain}

Let $W_0$ be the one-handle coefficient trace, $W_1$ its quotient by all
two-handle relations, and $W_2$ the quotient of $W_1$ by the three-handle
relations.  Denote the selected class in $W_0$ by $x_0$.

\begin{lemma}[A linear detector survives a quotient]\label{lem:quotient-detector}
Let $V$ be a vector space, $R\subseteq V$ a subspace, and
$\lambda:V\to\Q$ a linear map with $R\subseteq\ker\lambda$.  Then there is a
unique $\bar\lambda:V/R\to\Q$ such that
$\bar\lambda\circ\pi=\lambda$.  If $\lambda(x)\ne0$, then $\pi(x)\ne0$.
\end{lemma}

\begin{proposition}[Abstract nonzero class after quotients]\label{prop:nonzero-W2}
Suppose Hypotheses~\ref{hyp:P1} and~\ref{hyp:P2} hold.  Then there is a
linear functional $\ell_2:W_2\to\Q$ such that
$\ell_2(q_{12}q_{01}x_0)=2624$.  In particular, $q_{12}q_{01}x_0\ne0$.
\end{proposition}

\begin{proof}
Lemma~\ref{lem:HH0-descent} descends the quantum shadow detector through the
coefficient trace.  Hypothesis~\ref{hyp:P1} (status:
Proposition~\ref{thm:Cdischarge}) and Lemma~\ref{lem:quotient-detector}
descend it through $q_{01}$.  Hypothesis~\ref{hyp:P2} (status:
Proposition~\ref{thm:Sdischarge}) and the same quotient lemma descend it
through $q_{12}$.
At every stage the pullback agrees with the preceding row, so
Proposition~\ref{prop:D3-vT} gives the value $2624$.
\end{proof}

\begin{proposition}[Closed candidate class]\label{prop:closed-class}
Suppose Hypotheses~\ref{hyp:P2} and~\ref{hyp:P3} hold and an inhabitant of
Lean \texttt{ExternalGeometry} transports the class through the four-handle
isomorphism.  Then the class of Proposition~\ref{prop:nonzero-W2} maps to a
nonzero element of $\Stwo(X_J)_{0,494}$.  No such inhabitant is constructed.
\end{proposition}

\begin{proof}
MWW Theorems~3.7 and~3.10 identify the three-handle quotient, and
Proposition~3.4 identifies the four-handle stage once cited.
Hypothesis~\ref{hyp:P3} separates the computed empty-link ranks from those
citations.  The Lean implication in Theorem~\ref{thm:A} is conditional on
\texttt{ExternalGeometry}.
\end{proof}

\begin{remark}[Presentation and serialization dependence]\label{rem:serialization}
The raw coefficient $[h^3]\ell(\rho_h(W)-I)u$ depends on the based,
labelled, oriented movie presentation.  Reordering the same crossing rows
by the emitter's increasing-coordinate order produces a different series.
With the corrected collar endpoints its cubic coefficient is $0$ and its first
nonzero coefficient is $663872h^4$.  That order reverses $168$ non-disjoint
events and is not a legal serialization of the Johnson collar.  The former
value $-58976$ used the same mixed endpoint coordinates as the superseded
public value.  Presentation independence must be proved through simultaneous
transport of every part of the naturality square; see
Appendix~\ref{app:legality}.
\end{remark}

\section{A fully worked finite-dimensional example}\label{sec:toy-example}

This section illustrates the algebraic mechanism on a small example.  It is
not a four-manifold calculation and supplies no evidence for any
candidate-specific assumption.

Let $R_0=\Q[\eps]/(\eps^4)$, $t=1+\eps$, and $E_0=R_0^3$ with basis
$e_0,e_1,e_2$.  The two positive Burau generators are
\begin{align}
B_1&=
\begin{pmatrix}
-\eps&1+\eps&0\\
1&0&0\\
0&0&1
\end{pmatrix},
&
B_2&=
\begin{pmatrix}
1&0&0\\
0&-\eps&1+\eps\\
0&1&0
\end{pmatrix}.
\end{align}
Take the pure commutator word
$w_0=\sigma_1^2\sigma_2^2\sigma_1^{-2}\sigma_2^{-2}$.
Multiplying the eight matrices from left to right and reducing modulo
$\eps^4$ gives
\[
\rho(w_0)=
\begin{pmatrix}
1-\eps^3&-\eps^2+\eps^3&\eps^2\\
\eps^2&1&-\eps^2\\
-\eps^2+2\eps^3&\eps^2-3\eps^3&1+\eps^3
\end{pmatrix}.
\]
Let $u_0=e_0-e_1$ and $\ell_0=e_0^*$.  Define
$d_2(v)=[\eps^2]\,\ell_0(\rho(w_0)-I)v$.
Then $d_2=\begin{pmatrix}0&-1&1\end{pmatrix}$ and $d_2(u_0)=1$.
Model two relation maps by the columns
$r_0=(1,0,0)^T$ and $r_1=(0,1,1)^T$.  Then $d_2R=0$, so $d_2$ descends to
$E_0/\im R$ and the selected class is nonzero in the quotient.  What this
example lacks is precisely the geometric content: there is no claim that the
toy relation columns arise from handle maps or that the final quotient is a
smooth invariant.


\section{Lasagna modules and the comparison framework}

We recall only the features needed for the comparison.  For a smooth
compact oriented four-manifold $W$ and a framed link $L\subset\partial W$,
the skein lasagna module $\SN(W;L)$ is generated by decorated surfaces in
$W$, modulo local relations coming from $\KhR_N$ cobordism maps
\cite{MorrisonWalkerWedrich2019,ManolescuWalkerWedrich2023}.

For a handle decomposition
\[
W_1\subset W_2\subset W_3\subset W_4,
\]
where $W_1$ is a one-handlebody, $W_2$ adds two-handles along a framed link,
$W_3$ adds three-handles along embedded spheres, and $W_4$ adds four-handles,
the published results have the following roles.
\begin{itemize}
\item MWW Theorem 3.2 identifies the two-handle stage with a cabled skein
lasagna quotient, whose relations include beta and psi operations.
\item MWW Theorems 3.7 and 3.10 describe each three-handle as a coequalizer
of the two hemisphere maps, and give the complete handle formula.
\item MWW Proposition 3.4 says that a four-handle attachment induces an
isomorphism for the empty boundary link.
\item MWW Corollary 3.5 gives
$\SN(S^4)\cong\Z$, concentrated in bidegree $(0,0)$.
\end{itemize}

At $N=2$, a nonzero class in bidegree $(0,494)$ would therefore obstruct a
diffeomorphism to $S^4$.  The intended proof constructs a functional at the
one-handle stage, shows that it annihilates the two-handle and three-handle
relations, and transports the resulting nonzero class through the
four-handle isomorphism.

The formal algebra correctly captures the last sentence.  The next sections
formulate the geometric identifications required for P0, C and S and record
which parts are established and which remain \Open.

\section{Johnson product-ribbon presentation and P0}\label{sec:johnson-presentation}

We record the AR-side construction, summarize rejected compact lifts, and
give the Johnson-generator replacement for Hypothesis~\ref{hyp:P0}; its
remaining obligations are listed in Proposition~\ref{thm:P0discharge}.

\subsection{The actual AR link}

Use the coordinate-spine genus-three Heegaard splitting of $T^3$ from
Aitchison--Rubinstein \cite[pp.~5--12]{AitchisonRubinstein1984}; the complete
volume is available in a
\href{https://math.berkeley.edu/~kirby/papers/Gordon%20and%20Kirby%20%28editors%29%20-%20Four-manifold%20theory%20%28Durham%29%20-%20MR0780574.pdf}
{public Berkeley scan}.  Let
$C_i$ be the three coordinate spine circles with common base vertex $Q$, and
let $A_i$ be pairwise-disjoint narrow disk neighborhoods away from $Q$.
Write $t$ for the mapping-torus one-handle.  If $\psi_A$ is the
handlebody-preserving representative supplied by their Lemma~3.1, the core
of the $i$th mapping-torus two-handle is the embedded circle
\begin{equation}\label{eq:ARcore}
C_i^-\cup\lambda_i\cup\psi_A(C_i)^+\cup\mu_i,
\end{equation}
where $\lambda_i,\mu_i$ are the two arcs through $t$.  The framing annulus is
constructed simultaneously as
\begin{equation}\label{eq:ARannulus}
A_i^-\cup(\lambda_i\times I)\cup\psi_A(A_i)^+
       \cup(\mu_i\times I).
\end{equation}
The disks, cones and product levels in the AR construction are pairwise
disjoint, so \eqref{eq:ARcore}--\eqref{eq:ARannulus} define the whole
embedded framed link, not merely its words.

The remaining three fiber two-handles are the standard dual two-cells of the
coordinate-spine splitting.  After choosing orientations, denote their
boundary curves by
\[
r_{xy}=[x,y],\qquad r_{yz}=[y,z],\qquad r_{zx}=[z,x].
\]
Their embeddings are defined by the dual cells and their framings are the
fiber product framings.  The untwisted section-surgery handle $h_{CS}$ goes
geometrically once over $t$.  Thus the starting handle counts are
$(1,4,7,3,1)$ in indices $0,\ldots,4$: AR p.~8 removes one $3$- and the
$4$-handle from the mapping torus and glues
$S^2\times D^2=H^2\cup H^4$ in their place.

For an independently checkable standard-form bridge, set
\[
C_{AR}=\begin{pmatrix}0&0&1\\713&83&0\\-902&-105&-10\end{pmatrix},
\qquad
R=\begin{pmatrix}71&-466&1\\-610&4004&-7\\1&-7&-2\end{pmatrix}.
\]
Direct multiplication gives
\[
\det R=1,\qquad C_{AR}R=RA,
\]
and both matrices have determinant one and determinant-one difference from
the identity.  Hence the fiber diffeomorphism induced by $R$ transports the
complete AR construction to the displayed matrix $A$.  Since
$R\in GL^+(3,\mathbb R)$ and this group is connected, its derivative preserves
the homotopy class of the canonical normal framing of the section circle; in
particular, the untwisted choice $\varepsilon=0$ is preserved.

\subsection{Return to the linear torus map}

Let $\phi_A$ be the linear torus map, straightened to the identity on the
section ball, and choose an isotopy $\rho_u$, $u\in[-1,1]$, from $\phi_A$ to
$\psi_A$, stationary near its endpoints and relative to that ball.  Put
$g_u=\rho_u\phi_A^{-1}$.  With mapping-torus convention
$(x,-1)\sim(\phi(x),1)$, the formula
\begin{equation}\label{eq:mappingtorusF}
F([x,u]_{\phi_A})=[g_u(x),u]_{\psi_A}
\end{equation}
defines a diffeomorphism.  Indeed, $g_{-1}=\mathrm{id}$,
$g_1=\psi_A\phi_A^{-1}$, and therefore
$\psi_Ag_{-1}=g_1\phi_A$, which is precisely well-definedness at the seam.
Endpoint stationarity makes \eqref{eq:mappingtorusF} smooth, and the family
$g_u^{-1}$ gives its inverse.  It fixes the section ball.  Pulling the entire
AR handle decomposition back by $F$ replaces the top cores by the linear
images $\phi_A(C_i)$ and transports every component and framing annulus
simultaneously.

\subsection{Words and normal fields}

Lift $C_i$ and $\phi_A(C_i)$ to the universal cover.  The three top lifts
begin at the same point $A(Q)$.  Translate the whole top fiber by $-A(Q)$ and
extend this translation over its collar.  It is isotopic to the identity and
transports the entire link and its normal fields.  Since each attaching core
is unbased, changing the first face crossing only cyclically permutes its
word.  We may therefore take the top centerline to be the straight segment
from $0$ to $Ae_i$; it crosses the $j$th integer faces at
the rational times $k/(Ae_i)_j$.  A fixed arbitrarily small product
perturbation orders simultaneous events.  Reading the two base arcs and the
oppositely oriented bottom core gives exactly
\begin{equation}\label{eq:actualword}
m_i=t\phi_A(x_i)t^{-1}x_i^{-1}.
\end{equation}
This is the full event rule used by the public generator.

The framing is equally explicit.  AR push the coordinate axes in the
universal cover in direction $(1,1,1)$; the resulting strips have parallel
boundary components, and linearity of $A$ preserves this property on their
images \cite[pp.~16--17]{AitchisonRubinstein1984}.  Together with the
$\lambda_i/\mu_i$ rectangles, these strips are the normal fields in
\eqref{eq:ARannulus}.  Thus no blackboard integer is substituted for a
framing.

\subsection{The two cancellations}

Aitchison--Rubinstein identify $(t,h_{CS})$ as a complementary geometric
one/two-handle pair \cite[p.~7]{AitchisonRubinstein1984}.  Slide every other
two-handle passage over $h_{CS}$ along its product rectangle and cancel; this
removes the $t,t^{-1}$ passages with zero relative twist.

Now $Ae_1=e_3$, so the actual curve $m_1$ becomes $zx^{-1}$ and has one
geometric passage over the $x$ one-handle.  Thus $(x,m_1)$ is another
complementary pair.  Sliding every other $x$ passage over $m_1$ along parallel
product bands replaces $x$ by $z$ and cancels the pair.  All labels and
normal fields move with the bands.  The resulting counts are $(1,2,5,3,1)$,
and the exact word projection is
\[
\begin{aligned}
|m_2|&=311, & [m_2]_{(y,z)}&=(40,269),\\
|m_3|&=1460, & [m_3]_{(y,z)}&=(189,1271).
\end{aligned}
\]
For the committed compact $m_2$ word, a second linear and cyclic free
reduction still has length $311$; there is no adjacent inverse pair.  The
explicit 19-step Nielsen representative constructed in Appendix~B instead
reduces to length $309$.  It contains 40 $y/Y$ passages, rather than 42; with
the two $r_{xy}$ passages its detector has 42 channels rather than 44.  Thus
the 309/311 discrepancy is a genuine representative discrepancy, not merely
a printing convention.
The transported $r_{zx}$ has empty free word, but no split-unknot conclusion
is needed or inferred from that fact.

\subsection{The actual detector ball}

Cut the remaining genus-two handlebody along its $y,z$ disk system.  The
selected $m_2/r_{xy}$ state has 42 $y$ passages from $m_2$ and two from
$r_{xy}$.  At distinct product-normal levels, a regular neighborhood of these
44 passages is a standard three-ball $B$ containing a punctured-disk collar.
The six affine sweeps specify a pure 44-strand braid in this collar.  The
mapping-class interpretation of the braid group realizes it by an isotopy of
the punctured disk, and isotopy extension gives an ambient isotopy of $B$.
Choose the extension to be the identity to first order near every returned
puncture at its terminal time.  Hence each product normal returns, not merely
the total writhe.  The motion acts on every oriented normal copy
simultaneously.  Reconstruction of the $252$ crossing factors gives
$|B_{44}|=11340$ with $\operatorname{perm}(B_{44})=1$ and
$\operatorname{wr}(B_{44})=0$.

\begin{remark}[Rejected compact lifts]
Earlier symbolic lifts, including a 19-step Nielsen representative and
bounded IA corrections, fail channel count, free-basis, or embedded-collar
tests.  Appendix~\ref{app:retired-routes} records the falsified routes.
\end{remark}

\subsection{Johnson replacement}

Johnson's splitting-preserving generators for the genus-three splitting of
$T^3$ \cite{Johnson2011} give a more appropriate search space.  His generator
$\alpha_{ij}$ pushes the diagonal of the $x_i$--$x_j$ square to a
neighborhood of the spine; pushing to the opposite side gives a second lift of
the same transvection.  Factoring $A$ into $93$ unit transvections and choosing
a $93$-bit side pattern yields
\[
|m_2|=311,\qquad \#(y/Y)=42,\qquad p_y=42+2=44,
\]
with the three spine images forming a free basis of $F_3$ and the reduced
$m_2$ word agreeing with the compact word.  The class-two $r_{yz}$ coefficient
relative to the compact word is zero.

For each square move we record the diagonal, a two-edge path, and a square
normal.  Truncating both paths at one-quarter and three-quarters makes the
movie relative to the spine vertex.  If $M$ ranges over the $93$ intermediate
unimodular bases, the uniform radius
\[
r=\frac{1}{8\max_M\lVert M^{-1}\rVert_\infty}=\frac1{196104}
\]
is fixed throughout.

In the standard product chart $D^2\times[-1,1]$ for the reduced $y$
one-handle, take $B=D^2\times[-1/2,1/2]$.  The Johnson $m_2$ word has $41$
positive $y$ passages and one negative passage; $r_{xy}=[z,y]$ contributes one
positive and one negative passage.  Distinct rational points in $D^2$ give
pairwise-disjoint product arcs with constant transverse framing.  The oriented
dual rectangle for $r_{xy}$ has six point-push legs; expanding the resulting
$252$ geometric L/R factors gives an Artin word of length $11340$, which
agrees with the independently regenerated public word.

\begin{lemma}[Railroad linking]\label{lem:P0d-link}
The railroad reduced planar diagram of the Johnson CS handle picture satisfies
$\operatorname{lk}(m_2,r_{yz})=0$.
\end{lemma}

\begin{proof}
The mixed $m_2$--$r_{yz}$ signed crossing sum in the labelled reduced PD is
$0$, so the linking number vanishes.  By the Kirby handle-slide formula
\cite{Kirby1978}, this is the framing condition required for the product
cancellations above.
\end{proof}

\subsection{The three geometric identifications}

\begin{lemma}[Johnson--AR handlebody bridge]\label{lem:P0a}
The affine map
\[
S:\mathbb R^3/\mathbb Z^3\longrightarrow
\mathbb R^3/(2\mathbb Z)^3,\qquad
S(u)=2u-(1/2,1/2,1/2),
\]
is an orientation-preserving diffeomorphism carrying Johnson's standard
genus-three Heegaard pair to the coordinate-spine pair used by
Aitchison--Rubinstein.
\end{lemma}

\begin{proof}
Scale Johnson's period-one torus by $4$ and the AR period-two torus by $2$;
both become the period-four Freudenthal mesh $T$ ($64$ vertices, $384$
tetrahedra), on which $S(u)=2u-(1/2,1/2,1/2)$ is the translation
$T(v)=v-(1,1,1)$, a simplicial automorphism of $T$.  The Johnson spines
$K_J^0,K_J^1$ are the axis circles through $(0,0,0)$ and $(2,2,2)$, the AR
spines $K_{AR}^0,K_{AR}^1$ those through $(-1,-1,-1)$ and $(1,1,1)$; each is
a rose of three circles with $10$ vertices and $12$ edges (rank $3$).  In the
first barycentric subdivision $T'$ ($1664$ vertices, $9216$ tetrahedra) let
$H^0=N(K^0;T')=\bigcup_{\sigma\in K^0}D(\sigma)$ be the dual-block
neighbourhood of $K^0$ and $H^1$ the closure of its complement.  The
repository certificate exhibits explicit elementary-collapse sequences
$H^0\searrow K^0$ ($3944$ free-face collapses on $1440$ tetrahedra) and
$H^1\searrow K^1$ ($17688$ collapses on $7776$ tetrahedra) for both the
Johnson and the AR pair, each replayed step by step by an independent
checker; by Rourke--Sanderson, a compact $3$-manifold collapsing onto a
rank-three graph is a regular neighbourhood of it, hence a genus-three
handlebody.  The two handlebodies of each pair are disjoint, their union is
all of $T'$, and they share one closed connected boundary surface of $1056$
triangles and Euler characteristic $-4$.  The translation $T$ carries
$H_J^0$ onto $H_{AR}^0$ and $H_J^1$ onto $H_{AR}^1$ tetrahedron by
tetrahedron, while other shifts fail.  Regina~7.4 independently reports
all four complexes as genus-three handlebodies
(\texttt{recogniseHandlebody()}$=3$, $H_1=\mathbb Z^3$ on the simplified
triangulations, seven tetrahedra each).  The protected
ball of radius $1/196104$ about $0$ maps to the ball of radius $2/196104$
about $Q$.  No uniqueness of regular neighbourhoods is used.  Replay:
\texttt{scripts/verify\_elementary\_collapse.py --check} and
\texttt{scripts/verify\_t73\_handlebody\_bridge\_regina.py --check}.
\end{proof}

\begin{lemma}[The two framed cancellations]\label{lem:P0b}
In the Johnson--AR replacement, the pairs \((t,h_{CS})\) and \((x,m_1)\)
satisfy the geometric Kirby \(1/2\)-cancellation criterion.  The
cancellations transport the whole labelled attaching link and preserve the
section framing class \(\varepsilon=0\).
\end{lemma}

\begin{remark}[Status of Lemma~\ref{lem:P0b}]
\Open.  Established: for the selected Johnson lift, $\psi_A(x)=z$ at the
level of free-group words, so $m_1=zx^{-1}$ after the first cancellation,
and the certificate contains two standard local cancellation movies in
$\mathbb R^3$ with geometric intersection $1$.  Not established: the actual
PL homeomorphism $\psi_A$ composed from the $93$ Johnson generators, the
actual framed attaching link \eqref{eq:ARcore}--\eqref{eq:ARannulus} built
from it, and the cancellations of $(t,h_{CS})$ and $(x,m_1)$ performed on
that link with triangulated belt spheres, band slides, transported
components and transported normal fields.  Word-level identities do not
substitute for these embedded objects.
\end{remark}

\begin{lemma}[Compatibility with MWW cabling]\label{lem:P0c}
The framing on the 44 Johnson lanes is the framing used by the MWW cable
construction.
\end{lemma}

\begin{remark}[Status of Lemma~\ref{lem:P0c}]
\Open.  Established: $44$ height-monotone polylines with constant product
normals in a triangulated $3$-ball, labelled by the Johnson $y$-wickets,
whose derived crossing movie recovers the public $11340$-letter word.  Not
established: that these strands are the strands cut from the actual
cancelled attaching link by the $y/z$ disk system, and that their normals
are the MWW cable framing.  The committed strands are generated from the
word by a control construction; they calibrate the crossing extractor and
do not bind the framing.
\end{remark}

\subsection{P0 conclusion}\label{sec:p0-conclusion}

\begin{proposition}[Status of P0 for the Johnson replacement]\label{thm:P0discharge}
The replayed finite certificates establish: (i) the $93$-transvection
factorization of $A$ and the $93$-bit side pattern with a GAP free basis;
(ii) the exact compact $m_2$ word with $42$ $y$-passages and
$\operatorname{lk}(m_2,r_{yz})=0$ (Lemma~\ref{lem:P0d-link}); (iii) the
six-sweep word equal to the public $11340$-letter word; (iv) the
handlebody bridge of Lemma~\ref{lem:P0a}: dual-block regular
neighbourhoods of the Johnson and AR spines with explicit collapse
sequences, carried onto each other by $S$.  The
following obligations of Hypothesis~\ref{hyp:P0} are \Open:
(a) the identification of the mapping-torus handlebodies with this pair,
which the Aitchison--Rubinstein construction requires but the certificate
does not use; (b) the actual PL homeomorphism $\psi_A$ composed
from the $93$ Johnson generators and the actual Aitchison--Rubinstein framed
link \eqref{eq:ARcore}--\eqref{eq:ARannulus}; (c) the actual Kirby
cancellations of $(t,h_{CS})$ and $(x,m_1)$ with belt spheres, band slides
and transported normal fields (Lemma~\ref{lem:P0b}); (d) the collar strands
cut from the actual cancelled link and the framing identity with the MWW
cable (Lemma~\ref{lem:P0c}); (e) the identification of the closed picture
with $\Sigma_A^0$.  Hypothesis~\ref{hyp:P0} is therefore \Open.
\end{proposition}

\begin{proof}
Items (i)--(iv) are the replayed certificates listed in
Appendix~\ref{sec:reproducibility-extended}; item (iv) is the proof of
Lemma~\ref{lem:P0a}.  Items (a)--(e) are not established by them: the
committed collar strands are generated from the public word by a control
construction, and the cancellation movies are local models in
$\mathbb R^3$ rather than moves on the actual attaching link.
\end{proof}

\section{The coefficient-trace comparison}\label{sec:c-comparison}

This section records the C argument for the Johnson replacement collar.

\subsection{The product Hattori equivalence}

The compact words give the letter counts
\[
p_y=42+2=44,\qquad p_z=269+2=271.
\]
In the cyclic $m_2$ word, every one of the 42
$y/Y$ events is followed immediately by a distinct $z/Z$ event; the same
holds for the two $y/Y$ events of $r_{xy}$.  The witness lists all 44 pairs.
Those pairs are realized as 44 product ribbons of the P0 control
strands along their recorded normals, with $227$ model $z$-circles placed
off the P0 cube.  The scope is the Johnson replacement collar, not an
embedded cut in $\partial W_2$.

Let $\mathcal C$ be the pregraded rational tangle category with objects
$T:P_{86}\to P_{88}$, and let $F$ be the framed west-to-east tangle formed
by those rectangles.  MWW's one-handle theorem identifies the cut coefficient
as
\[
M_R(T,T')=\KhR_2(R\cup T'\cup\overline T)
\]
with its displayed shift and gluing actions
\cite[Theorem~4.7]{ManolescuWalkerWedrich2023}.  Pivotal tangle duality and
K\"unneth over $\Q$ give graded isomorphisms on this replacement collar
\begin{equation}\label{eq:HattoriActual}
H_{T,T'}:M_R(T,T')\xrightarrow{\cong}
\operatorname{Hom}_{\mathcal C}(FT,FT')\{-44\}
\otimes A^{\otimes227},
\qquad A=\Q[X]/(X^2).
\end{equation}
These maps form a coefficient-bimodule isomorphism on the replacement collar.
Gluing $f:S\to T$ before the source link and then applying the product
isotopy is isotopic rel boundary to first applying the isotopy and
precomposing by $Ff$.  The right action is the same argument with
postcomposition.  Thus both action squares are associativity of gluing for
one simultaneous surface.  At the chain level the two squares are
\[
\begin{array}{ccc}
C_R(T,T')&\xrightarrow{\ H_{T,T'}\ }&
C(FT,FT')\otimes C(S^1)^{\otimes227}\\
\big\downarrow {L_f}&&\big\downarrow {L_{Ff}\otimes1}\\
C_R(S,T')&\xrightarrow{\ H_{S,T'}\ }&
C(FS,FT')\otimes C(S^1)^{\otimes227},
\end{array}
\qquad
\begin{array}{ccc}
C_R(T,T')&\xrightarrow{\ H_{T,T'}\ }&
C(FT,FT')\otimes C(S^1)^{\otimes227}\\
\big\downarrow {R_g}&&\big\downarrow {R_{Fg}\otimes1}\\
C_R(T,U)&\xrightarrow{\ H_{T,U}\ }&
C(FT,FU)\otimes C(S^1)^{\otimes227}.
\end{array}
\tag{17}
\]
Here $C_R$ and $C$ are the strict Blanchet--Khovanov foam complexes.  Both
composites in each square are the same glued foam after exchanging two
disjoint collars.  BHPW strict functoriality removes the projective sign, and
taking homology gives \eqref{eq:HattoriActual} and both MWW action squares
for the Johnson replacement collar.  Equation~\eqref{eq:HattoriActual} is
not a chain-level complex of an embedded cut in $\partial W_2$.

The ordinary quotient statement, including both cyclic-relation spans, is
checked in the formal development.\footnote{See
\texttt{Smooth4PC/RepresentableCoefficient.lean}.}

\begin{lemma}[Actual coefficient bimodule]\label{lem:C1}
For the P0 replacement, \eqref{eq:HattoriActual} is a natural isomorphism of
the actual MWW coefficient bimodule, with both gluing actions.  Moreover the
map \(\operatorname{Sh}_q\) below is the quantum trace of that bimodule, not a
dimension-matched substitute.
\end{lemma}

\begin{remark}[Status of Lemma~\ref{lem:C1}]
\Open.  Established: the $44$ model rectangles from the P0 control strands
with product-normal $z$-sides, giving the Johnson $y$-then-$z$ pairing and
$227$ model $z$-circles, and the representable $HH_0$ reduction of
Section~\ref{sec:exact-algebra}.  Not established: the isotopy
\[
R\cup\overline T\cup T'\;\cong\;\overline{FT}\circ FT'\sqcup 227\text{ circles}
\]
for the actual cut tangle $R$, fixing the boundary, the detector ball and
the gluing supports, and hence the chain-level action squares (17).  Without
that isotopy, \eqref{eq:HattoriActual} is a model isomorphism and the map
$\operatorname{Sh}_q$ below is not yet identified as the quantum trace of
the actual coefficient bimodule.
\end{remark}

\subsection{Quantum trace and the endpoint map}

Set $R_q=\Q[q,q^{-1}]$.  The pregraded cyclic relation is
\[
L_f(m)=q^{|f|}R_f(m).
\]
By Lemma~\ref{lem:C1}, \eqref{eq:HattoriActual} and the counits are homogeneous
coefficient morphisms on the Johnson replacement collar.  BPW's canonical vertical-to-horizontal functor and the BHPW strict
tangle/foam functor give
\begin{equation}\label{eq:actualShadow}
\operatorname{Sh}_q:q\operatorname{Tr}(\mathcal C;M_R)
\longrightarrow\operatorname{Hom}_{R_q}(E_{86},E_{88}),
\end{equation}
after the flat-base qHH$_0$/Chern identification
\cite{BeliakovaPutyraWehrli2019,BeliakovaHogancampPutyraWehrli2023}.  Here
\[
E_{86}=(V^{\otimes86})_{\lambda=86},\qquad
E_{88}=(V^{\otimes88})_{\lambda=86};
\]
the first has rank one and the second rank 88.  The tangle maps are
$U_q(\mathfrak{sl}_2)$ intertwiners before restriction to these weight
spaces.

Base change $q\mapsto1+h$ factors through rational localization and the
completion of a Noetherian local ring, hence is flat.  Reduction modulo $h$
of the trace cokernel simply replaces $q^{|f|}$ by one, giving the ordinary
MWW coefficient $HH_0$; right exactness of tensor product suffices for this
last assertion.

\subsection{Selected class and divided detector}

Let $U:P_{86}\to P_{88}$ be the selected one-cup tangle in the P0 endpoint
order and put $T_1=F^{-1}U$.  Define
\[
v_T=H_{T_1,T_1}^{-1}
  (\operatorname{Id}_U\otimes X^{\otimes227}).
\]
The circle counits give
$\operatorname{Sh}_h([v_T])=u_h$.  Its absolute degree is
\[
-44+227+315-4=494.
\]

The normalized checked R-matrix on adjacent one-defect basis vectors is
\[
\begin{pmatrix}1-q^{-2}&q^{-1}\\q^{-1}&0\end{pmatrix}.
\]
After the position basis change and $t=q^{-2}$, it is the public positive
Burau block; the negative block is its inverse.  The physical cable has mixed
orientations.  BPW's oriented-cabling formula identifies it with the
all-positive model up to a writhe-dependent grading shift and pivotal basis
maps.  The complete public word has writhe zero, so the global shift
vanishes.  Any degree-zero permutation or pivotal chart $P$ transports the
operator, cup and cap simultaneously:
\[
W\mapsto PWP^{-1},\qquad u\mapsto Pu,\qquad
\ell\mapsto\ell P^{-1}.
\]
The matrix coefficient is unchanged; we choose the public endpoint labels
after this common transport.  The constant terms are fixed by the single
collar endpoint convention of Section~\ref{sec:exact-algebra}: the selected
cup joins the physical endpoints bound to collar positions $2$ and $87$, and
BPW's cup/cap formulas give
\[
u_0=e_2-e_{87},\qquad
\ell_0=e_{87}^*-e_2^*,
\]
as in \eqref{eq:u} and \eqref{eq:ell}, with no undetermined sign.  The
single authority for these conventions is the repository file
\texttt{data/T73\_ENDPOINT\_CONVENTION.json}, which records for each of the
$88$ endpoints its physical identity, owner, orientation, geometric order
(the frozen MWW cut slot), public order (the collar position), THXY index
and pivotal coefficient $\pm q^k$.  The program
\texttt{scripts/build\_t73\_endpoint\_transport.py} derives the monomial
matrix $P(q)$ from it, verifies
$W_{\mathrm{public}}=PW_{\mathrm{geometric}}P^{-1}$ letter by letter along
the actual $45{,}360$-letter cabled word with strand tracking (all eight
orientation patterns occur), and derives $u_0$ and $\ell_0$ from the
canonical coevaluation and the pivotal evaluation instead of reading them;
\texttt{recompute\_t73\_delta3.py} consumes that derivation.  All
omitted factors are invertible $q$-monomials or positive-order
$h$-corrections, and by Lemma~\ref{lem:filtered-cubic} below they do not
affect the divided cubic.

Over $\Q[[h]]$ define
\[
D_h=\ell_h(\rho_h(W)-I)\operatorname{Sh}_h.
\]
This is well defined on the quantum coefficient trace because
\eqref{eq:actualShadow} has already imposed quantum cyclicity.  The exact
Magnus calculation and Darn\'e's Andreadakis theorem give
$W\in\Gamma_3(P_{44})$ \cite[Theorem~6.2]{Darne2021}.  Since every pure
braid generator acts as $I+O(h)$, Lemma~\ref{lem:commutator-order} gives the
one-directional implication
\[
W\in\Gamma_3(P_{44})\ \Longrightarrow\
\rho_h(W)-I\in h^3\operatorname{End}(E_{88});
\]
the converse is not asserted and is not needed.  Independently of the
Andreadakis theorem, direct evaluation of all $7{,}744$ entries of
$\rho_h(W)-I$ verifies the membership on the right.
Thus $D_h$ takes values in $h^3\Q[[h]]$.  Division by $h^3$ is unique, and
reduction modulo $h$ gives a lift-independent ordinary functional
\[
D_3:HH_0(\mathcal C;M_R)\longrightarrow\Q.
\]
Positive-order corrections in the cup, cap, and basis change do not alter
the cubic coefficient (Lemma~\ref{lem:filtered-cubic}).

\begin{lemma}[Filtered cubic pairing]\label{lem:filtered-cubic}
Let $A(h)=\sum_{j\ge0}A_jh^j$ be an endomorphism series of $E_{88}$ with
$A_0=A_1=A_2=0$, and let $u(h)=\sum_k u_kh^k$ and
$\ell(h)=\sum_i\ell_ih^i$ be a vector series and a row series.  Then
\[
[h^3]\,\ell(h)A(h)u(h)=\ell_0A_3u_0,
\]
and the coefficient is unchanged when $(A,u,\ell)$ is replaced by
$(PAP^{-1},Pu,\ell P^{-1})$ for a constant automorphism $P$.
\end{lemma}

\begin{proof}
$[h^3]\ell(h)A(h)u(h)=\sum_{i+j+k=3}\ell_iA_ju_k$, and every term with
$j\le2$ vanishes.  Under simultaneous transport each summand is unchanged.
Both statements are formalized in
\texttt{Smooth4PC/FilteredCubicNaturality.lean}
(\texttt{pairingCoeff\_three\_of\_startsAtThree},
\texttt{cubic\_invariant\_under\_simultaneous\_transport}); the numerical
hypothesis $\rho_h(W)-I\in h^3\operatorname{End}(E_{88})$ is recomputed on
all $7{,}744$ entries by the endpoint transport program, and the cubic
agrees whether it is evaluated with the constant terms, with the full
transported $q$-series, or as $\ell_0A_3u_0$.
\end{proof}

The four mixed-index Burau sign variants
\[
-53824,\quad -53760,\quad -59008,\quad -59072
\]
are nonzero, but they are not the frozen cubic: each mixes two endpoint
coordinate systems.  The frozen public receipt and Lean give
\begin{equation}\label{eq:selectedD3}
[h^3]\,\ell_h(\rho_h(W)-I)u_h=2624\ne0
\end{equation}
in the single collar convention.  Equation~\eqref{eq:selectedD3} is a
statement about the frozen Burau computation; it becomes
$D_3([v_T])=2624$ only when Lemma~\ref{lem:C1} identifies $[v_T]$ with an
actual coefficient-trace class, which is \Open.  Nonvanishing does not
close C or S.

\subsection{The complete two-handle cocone}

\begin{lemma}[All-cable constant term]\label{lem:C2}
For every finite cable level and every sign-preserving cable braid \(b\),
the endpoint operator has constant term equal to its signed permutation.
After standardizing that permutation, the remaining pure-braid operator is
\(I+O(h)\).
\end{lemma}

\begin{proof}
BPW's oriented-cabling action is obtained from the strict BHPW tangle
functor.  At \(q=1\), the braiding on the fundamental representation and its
dual is the ordinary symmetric interchange, including the pivotal
identifications for oppositely oriented factors.  Therefore the constant
term of a braid is exactly the map of its signed endpoint permutation.  A
pure sign-preserving braid has identity permutation and hence constant term
\(I\).  All matrix entries are Laurent polynomials in \(q\), regular at
\(q=1\); after \(q=1+h\), subtracting the constant term leaves a multiple of
\(h\).  The argument is independent of the number of cable copies and so
holds at every finite MWW level.  This statement concerns the endpoint
representation only and does not assert the unknown symmetric-group
factorization of the anti-parallel MWW beta action.
The lemma is applied to the Johnson replacement after Lemma~\ref{lem:C1}.
\end{proof}

Only $D_3$ is asserted to descend after C1.  The remainder of this
subsection is the C2 argument after Lemma~\ref{lem:C1}.

P0c supplies product normals on the reconstruction strands, hence disjoint
parallel levels in the replacement collar.  We first type the construction at every MWW
cable state.  In owner order $(m_2,m_3,r_{xy},r_{yz},r_{zx})$, put
\[
n_y=(42,189,2,2,0),\qquad n_z=(269,1271,2,2,0).
\]
For $r\in\mathbb N^5$, cutting the actual cable $K(r,r)$ along the $y$- and
$z$-cocores gives
\[
p_y(r)=\sum_i r_i n_{y,i},\qquad
p_z(r)=\sum_i r_i n_{z,i}.
\tag{24}
\]
By Lemma~\ref{lem:C1}, once established, parallel copies of the P0 rectangles would give a chain isomorphism $H_r$ of the
same form as (17), with the corresponding endpoint sets.  The
$r_{zx}$ component is not assumed split; its stabilization relations are
the subject of Section~\ref{sec:local-stabilization}.  Since all
copies lie in one fixed tubular product chart, $H_r$ commutes with every
gluing action and with adding a parallel pair.

Let $s_0=e_{m_2}+e_{r_{xy}}$.  If $r\ge s_0$, let $\Omega_r$ be the finite set
of choices of one negative and one positive physical copy of each of $m_2$
and $r_{xy}$.  For $\omega\in\Omega_r$, close every unselected product factor
by the MWW core counit, use the selected 44 passages for the point-push, and
write $D_{r,\omega}$ for the cubic coefficient of the resulting row.  Define
\[
D_r=|\Omega_r|^{-1}\sum_{\omega\in\Omega_r}D_{r,\omega}.
\tag{25}
\]
For a state not above $s_0$, add the missing pairs with the once-dotted MWW
maps and define $D_r$ by pulling (25) back.  Disjoint supports make the order
of these additions immaterial.  Thus $D_r$ is defined on
every summand in
MWW Definition~3.1, including the states below the selected one.

For beta, apply the P0 collar motion simultaneously to all physical copies.
This is the oriented-cabling action of
\cite[Theorem~E]{BeliakovaPutyraWehrli2019}, made strict by BHPW.  Constant
permutations reindex $\Omega_r$.  After a placement is standardized, the
remaining pure braid is \(I+O(h)\) by Lemma~\ref{lem:C2}.  Simultaneous transport of the
operator, cup and cap, together with $\rho(W)-I=O(h^3)$, changes the row only
in order at least four.  Therefore
\[
D_r\beta_i(b)=D_r
\tag{26}
\]
for every $b\in B_{r_i,r_i}$.  This cubic-order argument does not use the
assertion that anti-parallel braid actions factor through a symmetric group,
which MWW leave open
\cite[Remark~3.3]{ManolescuWalkerWedrich2023}.

For a pair addition, split the target placements according to whether the
new pair is selected.  A beta move exchanges a selected new pair with an old
one, so (26) reduces both subsets to the same local core calculation.  On the
whole source that calculation is
\[
(\epsilon\otimes\epsilon)\Delta(1)=0,\qquad
(\epsilon\otimes\epsilon)\Delta(X)=1.
\]
For $r_{zx}$ the same equations are required inside a tubular neighbourhood
of the component, without any split hypothesis; this is the local statement
of Section~\ref{sec:local-stabilization}.  Forgetting the distinguished new pair has constant-size fibers
between placement sets, and the normalizations in (25) cancel their
cardinalities.  The dotted raw degree $+2$ is cancelled by the target cabled
shift $-2$.  Consequently, for every owner and state,
\[
D_{r+e_i}\psi_i^{[0]}=0,\qquad
D_{r+e_i}\psi_i^{[1]}=D_r.
\tag{27}
\]
Below $s_0$ the second equality is the definition; the Frobenius identities
and disjoint-support interchange prove the first equality and the commutation
of all owner squares.  The algebra of this statewise construction is
formalized in \texttt{Smooth4PC/ReynoldsCableCocone.lean}: for every state
$r\in\mathbb N^n$, the Reynolds average is invariant under a braid action
that permutes placements, the two pair-addition relations follow from
placement-wise relations with constant fibres, and the extension of $D_r$
to the states below $s_0$ along iterated once-dotted additions is
independent of the common upper bound and satisfies the $\psi^{[1]}$
relation at every state.  Every geometric input of that file (that the
actual MWW beta action permutes placements, that the actual pair additions
forget a distinguished pair with constant fibres and commute, and the
descent above $s_0$) is a hypothesis there, and for the actual cabled
source these hypotheses are part of the \Open\ Hypothesis~\ref{hyp:P1}.

Once (26)--(27) are established on the actual cabled source, together with
the coefficient trace relation of (17), they prove on the whole cabled
direct sum
\[
D_3(\beta_i(b)x-x)=0,\quad
D_3(\psi_i^{[0]}x)=0,\quad
D_3(\psi_i^{[1]}x-x)=0,
\]
and the quotient universal property with MWW Theorem~3.2 then gives
$\overline D_3:\mathcal S^2_0(W_2;\Q)\to\Q$.  At present (17) and
(26)--(27) are verified only on the control collar, so this descent is part
of the \Open\ Hypothesis~\ref{hyp:P1}.

\subsection{Local stabilization without a split hypothesis}
\label{sec:local-stabilization}

For a framed component $K_i$ and a cable state $r$, let $\chi_{i,r}$ denote
the local closure of the $r_i$ parallel copies of $K_i$ inside a tubular
neighbourhood of $K_i$, defined in the pre-quotient foam category by the
undotted and once-dotted sphere evaluations.  The statement needed for the
$\psi$ relations of every owner, including $r_{zx}$, is
\[
\chi_{i,r+e_i}\circ\psi_i^{[0]}=0,\qquad
\chi_{i,r+e_i}\circ\psi_i^{[1]}=\chi_{i,r},
\]
with no assumption that $K_i$ is split from the rest of the link.  This
local stabilization--retraction statement is \Open.  The earlier text
inferred a split unknot from the empty free word of the transported
$r_{zx}$; that inference is withdrawn, since a free-group word does not
determine an embedded component.

\subsection{C conclusion and the degree-$494$ square}

The MWW one- and two-handle maps fit into the commuting diagram
\[
\begin{array}{ccccccc}
M_R(T_1,T_1)&\xrightarrow{H_{T_1,T_1}}&
\operatorname{End}(U)\{-44\}\otimes A^{\otimes227}
&\xrightarrow{\operatorname{id}\otimes\epsilon^{\otimes227}}&
\operatorname{End}(U)&\xrightarrow{d_3}&\Q\\
\big\downarrow{\mathrm{Glue}_{1h}}&&&&&&\big\Vert\\
\mathcal S^2_0(W_1;K(s_0,s_0))\{315\}
&\xrightarrow{\ \iota_{s_0}\{-4\}\ }&
\mathcal S^{2,0}_0(W_1;K)/(\beta,\psi)
&\xrightarrow[\cong]{\ \Phi\ }&\mathcal S^2_0(W_2)
&\xrightarrow{\overline D_3}&\Q.
\end{array}
\tag{28}
\]
The left vertical map and its shift are MWW Theorem~4.7, and the bottom
isomorphism is MWW Theorem~3.2.  The upper selected element is
$\operatorname{Id}_U\otimes X^{\otimes227}$.  Its successive degrees are
\[
-44+227=183,\qquad183+315=498,\qquad498-4=494.
\]
Under Hypothesis~\ref{hyp:P1}, diagram (28) sends it to a class $x_2$
with $\overline D_3(x_2)=2624\ne0$ on the Johnson replacement collar.

\begin{proposition}[Status of C]\label{thm:Cdischarge}
Hypothesis~\ref{hyp:P1} is \Open\ for the Johnson replacement.
Established: the representable $HH_0$ reduction, Lemma~\ref{lem:C2} for the
endpoint representation, Lemma~\ref{lem:filtered-cubic}, and the frozen
Burau value \eqref{eq:selectedD3}.  Not established: the actual
product-ribbon isotopy and chain-level action squares of
Lemma~\ref{lem:C1}; the statewise definition of $D_r$ by placement averages
and its relations (26)--(27) on every cable state of the actual source; the
local stabilization statement of Section~\ref{sec:local-stabilization}; and
hence the descent of $D_3$ through the complete two-handle cocone and the
nonvanishing of the selected degree-$494$ class after all MWW two-handle
relations.
\end{proposition}

\begin{proof}
The established items are the replayed certificates and the Lean files
named in Appendix~\ref{sec:reproducibility-extended}.  The remaining items
require the actual cut tangle, which is not part of the current
certificates.
\end{proof}

\paragraph{C/S firewall.}
The $B$-external naturality obtained in C concerns boundary cobordisms in
the coefficient comparison.  It does not identify the local action of an
essential sphere in $\partial W_2$ with an ordinary dotted-sphere evaluation;
that is the separate calculation in Lemma~\ref{lem:Sendpoint}.

\section{Three-handle sphere closure}\label{sec:s-closure}

This section records the intended S argument.  Closed-manifold HJ
Theorem~5.3 is not used to conclude that $B$ is fixed.  Remove the final $4$-handle and call the
remaining handlebody $W_3$.  Its boundary is $S^3$.  Reversing the three
$3$-handles attaches three three-dimensional $1$-handles to $S^3$, and hence
\[
\partial W_2\cong\#^3(S^1\times S^2).
\]
The actual attaching sphere system has connected complement, since sphere
surgery on it gives the connected boundary of $W_3$.

\begin{lemma}[Kernel invariance under changing a complete sphere system]
\label{lem:Ssystem}
Let \(\mathcal S\) and \(\mathcal A\) be complete sphere systems in
\(\#^3(S^1\times S^2)\).  If they are related by isotopy, permutation and
sphere slides, then the kernels of the two total MWW three-handle maps out of
\(\mathcal S^2_0(W_2)\) agree.
\end{lemma}

\begin{proof}
An isotopy of an attaching sphere gives an isomorphic handle attachment.
A permutation merely reorders disjoint three-handles.  A sphere slide is
exactly a \(3\)--\(3\) handle slide.  Standard handle calculus gives a
diffeomorphism between the resulting three-handle cobordisms which is the
identity on their incoming boundary \(W_2\).  Naturality of the skein-lasagna
maps therefore gives \(q_{\mathcal A}=\Phi q_{\mathcal S}\) for an
isomorphism \(\Phi\) of the target modules.  Hence their kernels agree.
\end{proof}

Choose a standard complete system
\(\mathcal A=(A_1,A_2,A_3)\) in the Johnson replacement reversed
$1$-handle picture, already disjoint from the P0 reconstruction cube $B$.
Closed-manifold HJ Theorem~5.3 is not used to conclude that $B$ is fixed.
It is used only with Lemma~\ref{lem:Ssystem}: any other complete system,
including a going-up Kirby attaching system, is related by isotopy,
permutation and slides in the closed manifold, so the kernels agree.
The current version of \cite{HorvatJablonowski2025} (arXiv v3) contains
Lemma~5.5 (spotted-ball slide lemma) and Lemma~5.7 (relative uniqueness of
complete sphere systems); its Theorem~5.3 is proved through Lemma~5.7,
which in turn uses Lemma~5.5.  This paper invokes only Theorem~5.3 in the
closed manifold.  It does not invoke Lemma~5.7 in its relative form, which
would be the statement needed to keep the detector ball $B$ fixed during
the sphere slides; that relative statement is not used and the
corresponding fixed-$B$ conclusion is not claimed.
In the reversed $1$-handle picture, three belt spheres miss the P0 cube,
with dual loops pairing to the identity, $b=0$ endpoint foams, and
disjointness from the C1 model link and C2 supports; these are properties
of the model picture, and the corresponding statements for the actual
$\partial W_2$ are \Open\ (Proposition~\ref{thm:Sdischarge}).

MWW's intrinsic reformulation of a $3$-handle quotient uses the local module
action of $\mathcal S^2_0(S^2\times D^2)$.  At $N=2$, let $A_0$ denote the
essential sphere with one dot and $A_1$ the undotted sphere.  The quotient
relations are
\[
A_0=1,\qquad A_1=0.
\]
If $J_j$ is the equator of $A_j$, MWW's K\"unneth identification gives the
following exact description of the two hemisphere maps:
\[
\begin{array}{ccc}
\mathcal S^2_0(W_2;J_j)&\xrightarrow{\ (\Psi_{\Delta^+},\Psi_{\Delta^-})\ }&
\mathcal S^2_0(W_2)\oplus\mathcal S^2_0(W_2)\\
\big\downarrow\wr&&\big\downarrow{\ell\oplus\ell}\\
\mathcal S^2_0(W_2)\otimes\Q\{1,X\}
&\xrightarrow{\ (\ell\circ(1\otimes\epsilon),\,\ell\circ\mu_{A_j})\ }&
\Q\oplus\Q.
\end{array}
\tag{30}
\]
Here $\epsilon(X)=1$, $\epsilon(1)=0$, while
$\mu_{A_j}(v\otimes X)=vA_0$ and
$\mu_{A_j}(v\otimes1)=vA_1$.  Therefore S is equivalent to the two
whole-source identities
\[
\ell(vA_0)=\ell(v),\qquad \ell(vA_1)=0
\quad\text{for every }v\text{ and }j=1,2,3.
\tag{31}
\]
This identifies the formal maps with the coequalizer pair in MWW
Theorem~3.7, but not their endpoint evaluations for the actual essential
spheres.

\begin{proposition}[Fixed-detector sphere replacement]\label{thm:Sgeometry}
The total three-handle kernel may be computed using the standard belt-sphere
system of the Johnson replacement reversed $1$-handle picture, which misses
$B$, while the detector in $B$ remains fixed.  For this system, S reduces to
the six equations (31).
\end{proposition}

\begin{proof}
The reversed $1$-handle picture of the Johnson replacement places three
belt spheres missing $B$, with dual loops whose handle segments meet the
owner belt sphere once and whose chart returns miss every belt cube.
Closed-manifold HJ Theorem~5.3 relates any other complete
system to this one by isotopy in the closed manifold; Lemma~\ref{lem:Ssystem}
identifies the kernels.  The detector remains in $B$, which those belt
spheres miss.  The six equations (31) are then Lemma~\ref{lem:Sendpoint}.
All of this concerns the model reversed $1$-handle picture; the
corresponding constructions in the actual $\partial W_2$ are \Open\
(Proposition~\ref{thm:Sdischarge}).
\end{proof}

\subsection{The essential-sphere endpoint comparison}

\begin{lemma}[Endpoint factorization for a standard sphere]\label{lem:Sendpoint}
For every \(j\), every cable state and every source element \(v\), the
constant term of the actual endpoint image of
\(\mu_{A_j}(v\otimes a)\) is the old detector image of \(v\) tensored with
the rank-two TQFT map of the punctured standard sphere.  Consequently the
divided cubic row satisfies
\[
\ell(vA_0)=\ell(v),\qquad \ell(vA_1)=0.
\]
\end{lemma}

\begin{proof}
Make \(A_j\) transverse to the two-handle cocores, and let \(b_j\) be the
number of core disks.  MWW's proof of Theorem~3.10 represents the
\(\Delta^-\) hemisphere, before the core disks are restored, by the connected
genus-zero surface obtained from \(A_j\) after deleting those disks and the
\(\Delta^+\) disk.  Choose a generic movie after cutting the one-handles.
Because \(A_j\) and the detector sheet are disjoint, every critical point of
the movie belongs entirely to one of the two surfaces.  Mixed events are
therefore only invertible Reidemeister/pivotal moves, represented at the
endpoint by cable braids; there are no mixed births, saddles or deaths.

At \(q=1\), Lemma~\ref{lem:C2} identifies every mixed braid with the
symmetry of the tensor category.  Naturality of births, saddles and deaths
with respect to that symmetry moves all mixed braids to the ends of the
movie.  The same endpoint permutation and pivotal chart transport the
detector row and the surface map, so they cancel by simultaneous
conjugation.  Thus the actual constant map is
\[
\operatorname{Id}_{\rm old}\otimes\Delta^{b_j-1}\quad(b_j>0),
\]
and is \(\operatorname{Id}_{\rm old}\otimes\epsilon\) when \(b_j=0\).
The restored two-handle core disks give the \(b_j\) actual counits.  Hence
\[
\epsilon^{\otimes b_j}\Delta^{b_j-1}(1)=0,\qquad
\epsilon^{\otimes b_j}\Delta^{b_j-1}(X)=1
\tag{32}
\]
for \(b_j>0\), with the same equations interpreted as
\(\epsilon(1)=0,\epsilon(X)=1\) when \(b_j=0\).
The \(1,X\) basis is normalized by the \(\Delta^+\) cap in MWW
Example~3.8, and BHPW strict functoriality prevents an independent sign in
the \(\Delta^-\) movie.

The full mixed braid maps have the form \(P(I+O(h))\).  The permutation
\(P\) has just been cancelled, and pre- or postcomposing a row beginning in
order \(h^3\) with \(I+O(h)\) changes only order \(h^4\) and higher.
Therefore the constant calculation (32) is exactly the calculation seen by
the divided cubic row.  In the model picture the three belt spheres miss
the P0 cube and the C1 model $z$-circles, so $b_j=0$ and the foams are the
counit of MWW Example~3.8.  For the actual $\partial W_2$ the spheres
$A_j$, the numbers $b_j$, the intersections with the two-handle cocores and
the hemisphere movies are not constructed, so the lemma is established for
the model picture only.
\end{proof}

The endpoint exchange square is the $b=0$ counit of
Lemma~\ref{lem:Sendpoint}.
\[
\begin{array}{ccc}
\mathcal S^2_0(W_2;J_j)&\longrightarrow&
\mathcal S^2_0(W_2)\oplus\mathcal S^2_0(W_2)\\
\downarrow&&\downarrow\\
\mathcal S^2_0(W_2)\otimes\Q\{1,X\}&\longrightarrow&\Q\oplus\Q,
\end{array}
\]
whose lower maps are the actual counit row and actual \(A_j\)-action by the
movie decomposition in Lemma~\ref{lem:Sendpoint}.

\begin{proposition}[Status of S]\label{thm:Sdischarge}
Hypothesis~\ref{hyp:P2} is \Open\ for the Johnson replacement.
Established: Lemma~\ref{lem:Ssystem} (kernel invariance under isotopy,
permutation and sphere slides) and the algebraic form (30)--(32) on the
model picture.  Not established: an actual complete sphere system
$A_1,A_2,A_3$ in the actual $\partial W_2$ with the conditions of HJ
Theorem~5.3 verified (embedded, pairwise disjoint, connected complement,
spherical homology basis, disjoint from the detector ball, simultaneous
surgery giving $S^3$); the actual equators, hemispheres, intersections with
the two-handle cocores, restored core disks and Morse movies; and the
identities $D(vA_0)=D(v)$, $D(vA_1)=0$ for all source elements and all
relative-class summands computed from them.
\end{proposition}

\begin{proof}
Lemma~\ref{lem:Ssystem} is proved above.  The remaining items depend on
Hypotheses~\ref{hyp:P0} and~\ref{hyp:P1} and on constructions that are not
part of the current certificates.
\end{proof}

\section{Final identifications}\label{sec:final-identifications}

MWW Theorem~3.10 identifies the iterated quotient with the module after the
three-handles, and Proposition~3.4 gives the four-handle isomorphism
\cite{ManolescuWalkerWedrich2023}.  Propositions~\ref{thm:P0discharge},
\ref{thm:Cdischarge}, and~\ref{thm:Sdischarge} record the status of
Hypotheses~\ref{hyp:P0}--\ref{hyp:P2}; all three are \Open.  For the Johnson
replacement four-handle picture $X_J$, attaching the three $3$-handles
along the belt spheres of the reversed $1$-handles restores the boundary
$S^3$ (Hypothesis~\ref{hyp:P3}), a PL $4$-ball is attached, and the quantum
degree-$494$ summand of the standard sphere module vanishes by the
computation in the status remark of Hypothesis~\ref{hyp:P3} together with
MWW Corollary~3.5.  The staged Kirby pipeline would identify $X_J$ with
$\Sigma_A^0$ once Hypothesis~\ref{hyp:P0} is established.

\begin{remark}[Formalization boundary]\label{rem:lean-boundary}
A Lean development \cite{Lean4} proves Theorem~\ref{thm:A} as a conditional
implication from structures \texttt{ExternalGeometry} and
\texttt{CSExternalGeometry}.  It checks the matrix arithmetic, the value
$2624$, the degree $494$, and the abstract quotient lemmas, but does not
construct the geometric instances proved in this paper.  Appendix~\ref{app:lean-fields}
lists the correspondence with the Lean fields.
\end{remark}

\begin{remark}[Earlier assumption register]\label{sec:retired-assumptions}
An earlier project draft grouped unresolved candidate geometry into five
assumptions A1--A5 (actual presentation, coefficient comparison, all-state
cocone, sphere maps, and rational closure).  The present interface is
Hypotheses~\ref{hyp:P0}--\ref{hyp:P3} together with the uninhabited Lean
fields \texttt{ExternalGeometry} and \texttt{CSExternalGeometry}.  None of
A1--A5 is used as a premise of Theorem~\ref{thm:joined}.  For the Johnson
replacement, A1 corresponds to P0 and P3/E13, A2--A3 to C, and A4 to S;
these are \Open.  The computed part of A5 is discharged by P3/E12 together
with the cited MWW Proposition~3.4 and Corollary~3.5.
\end{remark}


\section{Related work}

The Cappell--Shaneson literature contains both the original construction
\cite{CappellShaneson1976,AitchisonRubinstein1984} and substantial
standardness results
\cite{Akbulut2010,Gompf2010,KimYamada2017,Iwaki2024}.  Our use of Iwaki is
limited to the standard matrix form, the representative table, and the
homotopy-sphere criterion.  His table entry for $(41,189,73)$ is motivation,
not a singularity or exoticness theorem.

Manolescu--Neithalath and Manolescu--Walker--Wedrich provide the relevant
handle calculus for skein lasagna modules
\cite{ManolescuNeithalath2020,ManolescuWalkerWedrich2023}.  BPW quantum trace
functoriality and the strict BHPW tangle theory provide general comparison
tools \cite{BeliakovaPutyraWehrli2019,
BeliakovaHogancampPutyraWehrli2023}.  Section~7 applies these functors to the
actual replacement coefficient and proves the statewise naturality.

Ren and Willis use Khovanov-theoretic skein lasagna methods to distinguish
other exotic four-manifolds \cite{RenWillis2024}.  Their computations are
independent of the trace-73 construction and are not used as its comparison
theorem.

Horvat--Jab\l{}onowski's recent work gives a useful geometric criterion for
recognizing three-handle attaching systems \cite{HorvatJablonowski2025}.
Theorem~5.3 is used only with Lemma~\ref{lem:Ssystem} for kernel invariance
in the closed manifold, not to fix $B$.  Their Lemma~5.5 (spotted-ball slide
lemma) and Lemma~5.7 (relative uniqueness of complete sphere systems),
present in the current arXiv version, enter only through the proof of
Theorem~5.3 and are not invoked directly here.

\section{Conclusion}

For the Johnson replacement we have formulated
Hypotheses~\ref{hyp:P0}--\ref{hyp:P3} and recorded their status.  The
finite layer (the value $2624$ of the frozen Burau computation, the degree
$494$, and the third-order property of the public word) and the abstract
Lean implication are established.  Hypotheses~\ref{hyp:P0}--\ref{hyp:P2}
are \Open: they require the actual Cappell--Shaneson geometry, namely the
actual Johnson PL homeomorphism $\psi_A$, the actual Aitchison--Rubinstein
framed link, the actual Kirby cancellations and cut tangle, and the actual
sphere system with its hemisphere maps.  Together with Theorem~\ref{thm:A}
and Iwaki's homotopy-sphere criterion, this gives a conditional smooth
obstruction in quantum degree $494$ once those hypotheses are established
and the geometric data are assembled into the Lean interface
\texttt{ExternalGeometry}.  That interface is not constructed here, so no
counterexample is claimed.  Independent replay commands are listed in
Appendix~\ref{sec:reproducibility-extended}.

\appendix

\input{sec-retired-routes}

\section{P0 reconstruction protocol}\label{app:p0-reconstruction}

An admissible P0 certificate consists of explicit rational PL data for a
triangulated three-ball $B$, forty-four height-monotone labelled strands
$c_i:[0,1]\to B$ with normal vectors, passage binding to the reduced
attaching link, two local cancellation movies preserving framing, and an
ordered crossing movie.  The verifier checks passage binding, disjointness,
framing transport, and letter-by-letter equality of the induced
relative-endpoint braid with the public $11340$-letter word.

\begin{lemma}[Soundness of the P0 certificate]
Suppose an input passes the strict verifier, and suppose its independent
embeddedness, disjointness and local-movie receipts are valid for the stated
PL complex.  Then $B$ is an embedded framed three-ball in the reduced AR
handlebody, its $44$ strands form a framed collar, and their relative-endpoint
geometric braid equals the public $11340$-letter pure braid.
\end{lemma}

\begin{proof}
The PL ball and the passage-binding map place every $c_i$ in the actual AR
boundary chart and preserve ownership and endpoint order through the two
cancellations.  Height monotonicity makes the $c_i$ a braid in the collar;
the disjointness receipt makes the strands pairwise disjoint.  The normal
vectors and their transport receipts give the framing on each strand and
show that the two cancellation movies preserve it.  Reading the certified
crossing events in increasing height gives the relative-endpoint Artin word
of this embedded tangle.  The verifier's final equality identifies that word
letter-for-letter with the public word.  Since relative-endpoint isotopy
classes of labelled braids are represented by their Artin words, the claimed
braid equality follows.
\end{proof}

\begin{remark}
Symbolic witnesses that import crossing rows without embedded geometry fail
this contract.  The committed Johnson certificate satisfies the finite
contract, but its strands are a control realization of the public word
(Section~\ref{sec:johnson-presentation}), so its embeddedness and
disjointness receipts are receipts for that model, not for strands cut from
the actual attaching link.  The machine-checkable replay is listed in
Appendix~\ref{sec:reproducibility-extended}.
\end{remark}

\section{The complete degree-three endpoint vector}\label{app:vector}

With the conventions of Section~\ref{sec:exact-algebra}, the coefficient
vector $[\eps^3](\rho_t(B_{88})-I)u=(a_0,\ldots,a_{87})^T$ is
\begin{align}
(a_0,\ldots,a_{87})={}&(
-7052,-7052,164,164,332,332,324,324,316,316,308,308,\notag\\
&300,300,292,292,284,284,276,276,268,268,260,260,\notag\\
&252,252,244,244,236,236,228,228,220,220,212,212,\notag\\
&204,204,196,196,188,188,180,180,172,172,164,164,\notag\\
&156,156,148,148,140,140,132,132,124,124,116,116,\notag\\
&108,108,100,100,92,92,84,84,76,76,68,68,\notag\\
&60,60,52,52,44,44,36,36,28,28,20,20,\notag\\
&12,12,-164,-164).
\end{align}
The detector covector is $\ell=e_{87}^*-e_2^*$, so
$\ell(a_0,\ldots,a_{87})^T=a_{87}-a_2=-164-164=-328$.
The sum of the $88$ entries is zero, as expected for the unreduced Burau
difference on the invariant total-coordinate functional.

\section{The four grading contributions}\label{app:grading}

\begin{center}
\begin{tabular}{@{}p{0.17\textwidth}p{0.45\textwidth}p{0.27\textwidth}@{}}
\toprule
term & source & calculation\\
\midrule
$-44$ & remove the $p(N-1)$ Hom normalization for $p=44$, $N=2$ & $-44(2-1)=-44$\\
$+227$ & $227$ distinct Frobenius labels, each of degree $+1$ & $227\cdot1=227$\\
$+315$ & one-handle gluing shift from $44$ $y$-pairs and $271=44+227$ $z$-pairs & $44+271=315$\\
$-4$ & cabled shift for $N=2$, $\alpha=0$, $|r|=2$ & $(1-2)(2\cdot2+0)=-4$\\
\midrule
total & & $-44+227+315-4=494$\\
\bottomrule
\end{tabular}
\end{center}

\section{Convention-legality criteria}\label{app:legality}

A serialization is legal only if it describes the same framed handle
presentation, the same oriented based cut tangle, the same labelled boundary
pairing, the same normal field, and the same insertion point.  Its event order
must be a linear extension of the geometric partial order.  Adjacent events
may be exchanged only when their supports are disjoint and a strict
interchange theorem applies.  Word replacement requires equality in the
based labelled braid or tangle category; equality of closures, permutations,
free-group words, or homology classes is insufficient.

A coordinate change $P$ must act simultaneously:
\begin{equation}\label{eq:simultaneous-coordinate}
W'=PWP^{-1},\qquad u'=Pu,\qquad
\ell'=\ell P^{-1},
\end{equation}
with the shadow, collar, basepoint, and insertion point transported in the
same naturality square.  Under \eqref{eq:simultaneous-coordinate},
$\ell'(W'-I)u'=\ell(W-I)u$.  Thus coordinate invariance is proved when the
simultaneous transport exists; it is not imposed by selecting the preferred
numerical answer.

A collar is legal only if it is an actual framed tubular collar of the fixed
component and is isotopic to the registered collar relative to the boundary,
basepoints, labels, old input balls, normal framing, and insertion point.
These criteria were written before the detailed alternative collar outputs
were opened.  Their use in the Johnson replacement is discharged by P0.

\section{Terminology}\label{app:glossary}

\begin{description}
\item[Actual map]
A map induced by a specified geometric tangle or surface cobordism, as
opposed to an abstract linear map having the desired source and target.
\item[Candidate-specific binding]
An identification between immutable finite data and the actual geometric
object required by a published theorem.
\item[Complete quotient]
The quotient by all relation generators in the applicable MWW handle formula,
not by a sampled finite subset.
\item[Divided cubic]
The reduction modulo $h$ of a map known on its whole domain to be divisible by
$h^3$, as in \eqref{eq:D3}.
\item[Endpoint shadow]
The image of a coefficient trace class under the quantum
vertical-to-horizontal trace and the fixed endpoint functor.
\item[Finite certificate]
Immutable finite data plus a deterministic exact checker.  The term has no
geometric force beyond the predicates actually recomputed.
\item[Registered movie order]
The oriented point-push chronology fixed in the public primitive input.
\item[Johnson replacement]
The splitting-preserving $\alpha_{ij}$ lift, 93-bit side choice, and CS handle
picture discharged by P0--P3/E13 in this paper.
\end{description}

\section{Extended reproduction commands}\label{sec:reproducibility-extended}

The public repository is
\begin{center}
\url{https://github.com/toffee-desuwa/smooth4pc-t73-lean}.
\end{center}
The commands below assume a fresh clone and Python~3.10 or later.

\subsection{Reconstruct the detector}

\begin{verbatim}
python -I -B scripts/recompute_t73_delta3.py --check
\end{verbatim}
Expected bearing lines include
\texttt{B44\_LENGTH=11340},
\texttt{DELTA3\_ETA\_T1=2624},
\texttt{DELTA3\_XI=0}, and \texttt{VERIFY=PASS}.

\subsection{Replay Johnson P0, C, S, and P3 certificates}

\begin{verbatim}
python3 scripts/verify_elementary_collapse.py --check
python3 scripts/verify_t73_handlebody_bridge_regina.py --check
python3 scripts/build_t73_johnson_pl_generators.py --check
python3 scripts/build_t73_section_straightening.py --check
python3 scripts/verify_t73_pl_homeomorphism.py --all
python3 scripts/verify_t73_pl_homeomorphism.py --theta
python3 scripts/compose_t73_psi_A.py --check --crosscheck
python3 scripts/certify_t73_p0_johnson.py --check
python3 scripts/certify_t73_c1_cut_link.py --check
python3 scripts/certify_t73_c2_comparison.py --check
python3 scripts/certify_t73_s_relative_moves.py --check
python3 scripts/certify_t73_p3_four_handle.py --check
python -B scripts/certify_t73_e12_s4.py --check
python -B scripts/certify_t73_e13_close.py --check
python3 scripts/audit_t73_premises.py --check
python -B scripts/check_t73_claim_boundary.py
\end{verbatim}

\subsection{Compile and inspect the Lean development}

After installing \texttt{elan}, run
\begin{verbatim}
lake update
lake exe cache get
python -B tests/test_t73_minimal_formalization.py -v
lake lean T73Audit.lean
\end{verbatim}
The Python test compiles the audited modules into a fresh temporary output
root, requires exactly $38$ axiom reports, rejects any axiom outside
\texttt{propext}, \texttt{Classical.choice}, and \texttt{Quot.sound}, and
rejects \texttt{sorryAx}.

\subsection{Build this paper}

From \path{paper/spc4-t73-candidate}:
\begin{verbatim}
pdflatex -interaction=nonstopmode -halt-on-error main.tex
bibtex main
pdflatex -interaction=nonstopmode -halt-on-error main.tex
pdflatex -interaction=nonstopmode -halt-on-error main.tex
\end{verbatim}
The reviewed PDF is copied to
\path{output/pdf/spc4-t73-candidate.pdf}.


\section{Correspondence with the Lean fields}\label{app:lean-fields}

For readers checking the formal boundary, the principal allocation is:
\begin{center}
\small
\begin{tabularx}{\textwidth}{@{}l X@{}}
\toprule
Lean field & Mathematical obligation \\
\midrule
\texttt{x0}, \texttt{ell0}, \texttt{ell0\_x0} & P1 comparison and the selected genuine MWW class. \\
\texttt{q01}, \texttt{ell1\_comp\_q01} & Complete two-handle beta/psi quotient and descent. \\
\texttt{q12}, \texttt{ell2\_comp\_q12} & Reversed belt spheres of S, MWW hemisphere maps, and complete coequalizer.  Not the E7 ``$B$ is fixed'' claim. \\
\texttt{transport}, \texttt{fourIso} & E11 for the Johnson replacement $X_J$, not $\Sigma_A^0$. \\
\texttt{s4DegreeZero} & Computed \texttt{EmptyKhQ}(494)=0 and PL $S^4$ Euler data from \texttt{certify\_t73\_e12\_s4.py}.  The isomorphisms $G(B^4)\simeq G(S^4)$ and $\mathcal S^2_0(S^4)\cong\mathbb Z$ are MWW Proposition~3.4 and Corollary~3.5, cited not formalized.  Empty-link Lean inhabitant only. \\
\texttt{diffeomorphismEquiv} & Graded diffeomorphism invariance of lasagna modules. \\
\texttt{matrixConditionsToHomotopySphere} & E13 constructed CS handle picture; Lean \texttt{CSTopologyData} uninhabited. \\
\bottomrule
\end{tabularx}
\end{center}

The geometric inputs formulated in Sections~\ref{sec:johnson-presentation}--\ref{sec:final-identifications}
are themselves \Open, and they remain to be assembled into a single Lean
inhabitant of \texttt{ExternalGeometry}; that assembly is \Open\ for the Johnson candidate.
The empty-link control in \texttt{T73S4Inhabitant.lean} packages only
\texttt{EmptyKhQ}.  The Lean theorem consumes no stronger statement, and it
records the external topology as structure fields rather than formalizing it
internally.

\bibliographystyle{amsplain}
\bibliography{references}

\end{document}
